%% Nature Computational Science submission
%% Draft — [DATE]
\documentclass[pdflatex,sn-nature]{sn-jnl}

\usepackage{graphicx}
\usepackage{multirow}
\usepackage{amsmath,amssymb,amsfonts}
\usepackage{amsthm}
\usepackage[title]{appendix}
\usepackage{xcolor}
\usepackage{booktabs}
\usepackage{array}
\usepackage{tabularx}

\raggedbottom

\begin{document}

%%==============================================================%%
%% Title
%%==============================================================%%
\title[Temporal multimodal ALL prediction]{%
  Temporally-structured multimodal prediction across six clinical
  targets in paediatric acute lymphoblastic leukaemia}

%%==============================================================%%
%% Authors
%%==============================================================%%
\author*[1,2]{\fnm{xxx} \sur{xx}}\email{xxx}

\author[1]{\fnm{[Author 2 Given]} \sur{[Author 2 Family]}}

\author[1]{\fnm{[Author 3 Given]} \sur{[Author 3 Family]}}

\affil*[1]{\orgdiv{[Department]},
           \orgname{[HOSPITAL NAME / University]},
           \orgaddress{\city{[City]}, \country{China}}}

\affil[2]{\orgdiv{School of Computing and Information Systems},
           \orgname{The University of Melbourne},
           \orgaddress{\city{Melbourne}, \country{Australia}}}

%%==============================================================%%
%% Abstract
%%==============================================================%%
\abstract{%
Paediatric acute lymphoblastic leukaemia (ALL) is the most common
childhood malignancy, yet 15--20\% of patients relapse after
apparently successful treatment.
Existing risk-stratification approaches operate on single-timepoint
snapshots and have not been systematically evaluated across the full
clinical decision timeline nor validated across multiple independent
centres.
Here we present a temporally-structured multimodal prediction framework
evaluated on six clinical targets simultaneously: four classification
tasks (Final Risk, D19MRD, D46MRD, Relapse) and two survival endpoints
(event-free survival, EFS; overall survival, OS), whose feature sets
are strictly aligned with the clinical data-acquisition timeline from
hospital admission to Day~46.
An automated morphology pipeline converts bone marrow smear images
into structured cell-level profiles, and we systematically test—at
every feature level—whether raw images or derived morphology add
predictive signal beyond the clinical features—signal that current
risk-stratification rules leave unused.
Phase~1 evaluates four large language models (LLMs) in a zero-shot
setting; Phase~2 trains decision trees, random forests, and random
survival forests on a temporally partitioned case database.
In the internal cohort, Phase~2 achieves AUROC~0.949 for Final Risk
and concordance index~0.746 for EFS.
A cosine similarity retrieval-augmented generation (RAG) belief
revision improves D46MRD AUROC from 0.655 to 0.881 ($+0.226$) in
the largest training regime.
External validation on a merged 41-patient cohort pooled from two
hospitals confirms robust Final Risk transfer (LLM zero-shot AUROC up
to 0.925; Phase~2 RF up to 0.991) without retraining; pooling also
makes Relapse and D19MRD two-class and AUROC-evaluable, although the
cohort remains small and, because one source hospital is a
relapse-enriched referral series, external relapse metrics are
exploratory.
This framework demonstrates that temporally stratified multimodal
learning, aligned with clinical workflows, enables consistent
multi-target prediction that generalises to independent cohorts.
}

\keywords{paediatric acute lymphoblastic leukaemia, temporal feature
design, large language models, risk stratification, survival analysis,
retrieval-augmented generation, external validation}

\maketitle

%%==============================================================%%
\section{Introduction}\label{sec:intro}
%%==============================================================%%

Acute lymphoblastic leukaemia (ALL) is the most common childhood
malignancy, accounting for approximately 25\% of all paediatric
cancers~\cite{[CITE:Hunger2015]}.
Risk-adapted chemotherapy has substantially improved survival in
high-income settings, yet minority of patients who relapse face substantially worse outcomes, with long-term
survival dropping to 30--50\% upon re-induction~\cite{[CITE:Bhojwani2013],[CITE:Pui2008]}.
Because the majority of relapses occur after apparent clinical
remission, the prospective identification of high-risk patients—early
enough to intensify or de-escalate therapy—remains a central and
unresolved challenge in paediatric oncology. This challenge is
particularly important from a global-health perspective: contemporary
risk stratification increasingly depends on specialised laboratory
investigations whose cost, infrastructure requirements, and
availability differ substantially across health-care settings [Ref]. Thus,
the clinically optimal prediction strategy may not simply be the model
with the highest discrimination, but the model that achieves sufficient
predictive performance with the least additional diagnostic burden [Ref].

Current risk stratification in paediatric ALL relies
on pre-specified categorical cut-offs and combinations of age, white blood cell count,
cytogenetics, immunophenotype, and minimal residual disease (MRD) with the exact timing and criteria varying across treatment
protocols [Ref].
These rules do not fully exploit the multimodal information that
accumulates during treatment, nor do they yield personalised
probability estimates for future relapse or survival.
Computational approaches promise to fill this gap, but most published
models make simplifications that break down at the bedside: they treat
outcome prediction as a static, single-shot problem, when in reality
paediatric ALL management unfolds as a structured sequence of events in
time [Ref, a summary table in appendix and recap in discussion]. Moreover, model performance is rarely considered together with
the resources required to acquire its predictors. Age and routine
blood counts are inexpensive and widely accessible, whereas
multiparameter flow cytometry, cytogenetics, molecular fusion testing,
and repeated MRD measurements require progressively greater laboratory
capacity and expenditure [Ref].
Consequently, a modest increase in predictive accuracy from a
feature-rich model may entail a disproportionately large increase in
clinical acquisition cost [Ref from IEEE].

We argue that realistic clinical prediction is governed by two
distinct and largely independent timelines, and that existing methods
collapse or ignore all three.
The first is the \emph{feature-acquisition timeline}: a patient's data
do not arrive at once but in a fixed clinical order—white blood cell
count and age at admission (Day~0), immunophenotype (Day~1),
fusion genes, karyotype and FISH (Day~7--10), and
MRD (Day~19--46).
A model trained on features pooled across all timepoints implicitly
uses information unavailable at the moment of decision, producing
optimistic but non-deployable estimates~\cite{[CITE:Steyerberg2019]}.
The second is the \emph{prediction-target timeline}: clinicians do not
make one prediction but a sequence of decisions—early response
(D19MRD), end-of-induction response (D46MRD), Final Risk assignment,
and long-horizon outcomes (relapse, EFS, overall survival [OS]).
Studying any single target in isolation, as is common, obscures how one
data-acquisition workflow must support all six decision nodes.
Aligning a dataset and a method to both timelines simultaneously,
rather than one at a time, is the central novelty of this work and the
source of its clinical realism.

Building on this principle, we construct a dataset in which seven
temporally ordered clinical feature sets (six strictly nested;
$\mathcal{F}_4$ is a Day-1 branch of $\mathcal{F}_3$) map onto
the feature-acquisition timeline, six prediction targets map onto the
target timeline, and patient enrolment year defines the
data-accumulation timeline.
On this substrate we evaluate two complementary methods.
Phase~1 applies large language models (LLMs) in a zero-shot setting,
probing the predictive upper bound reachable with \emph{no} accumulated
training data—the regime a clinician faces when a new protocol or a new
hospital comes online~\cite{[CITE:Phaterpekar2026]}.
Phase~2 trains decision trees, random forests, and random survival
forests on the growing case database under temporal splits, emulating
prospective deployment as the data-accumulation timeline advances.
A retrieval-augmented generation (RAG) component then leverages the
accumulated database to apply a similarity-weighted belief revision to
model probabilities and to generate interpretable clinical narratives.
In parallel, an automated bone marrow morphology pipeline (SAM-based
cell segmentation and BiomedCLIP clustering) converts smear images
into a structured cell-morphology profile, letting us test—at every
level of the clinical feature hierarchy—whether imaging adds
prognostic information that current risk rules and prior computational
studies leave unused.

Our contributions are: (\textit{i})~a seven-level, three-timeline
clinical feature-and-target hierarchy aligned with the clinical
workflow; (\textit{ii})~Phase~1, a systematic zero-shot evaluation of
four LLMs across all feature levels and all six prediction targets,
establishing training-free baselines; (\textit{iii})~a systematic
test of whether bone marrow smear images—unused by current
risk-stratification rules—add predictive signal beyond clinical
features, finding that a structured morphology profile derived from
the images carries relapse-relevant signal at diagnosis; (\textit{iv})~Phase~2, a decision
tree, random forest, and random survival forest system trained on a
dynamic case database partitioned by enrolment year and validated
under three temporal splits; (\textit{v})~a cosine-similarity RAG
belief revision that refines model probabilities and generates
interpretable LLM clinical narratives; and (\textit{vi})~multi-site
external validation at independent centres that received no training
data.

%%==============================================================%%
\section{Results}\label{sec:results}
%%==============================================================%%

\subsection{Study cohorts and temporally-structured feature design}
\label{sec:cohort}

The internal training cohort comprised paediatric ALL patients treated
at [HOSPITAL NAME] between 2015 and 2019 (Table~\ref{tab:cohort};
$n=301$).
Prediction targets span the treatment timeline: Day~19 MRD positivity
(D19MRD $\geq 1\%$), Day~46 MRD positivity (D46MRD $\geq 0.01\%$),
final risk classification (Final Risk: LR vs IR/HR; 159 LR, 142 IR/HR),
and relapse occurrence (44 of 301, 14.6\%).
EFS events (relapse, induction failure, or death; 75 of 301, 24.9\%)
and OS events (death; 24 of 301, 8.0\%) together extend the prediction
targets to six.

External validation used a single pooled cohort of 41 patients
collected at two hospitals: Lianyungang First People's Hospital
(20 patients, diagnosed 2019--2023) and Children's Hospital of Fudan
University, Shanghai (21 patients, 2017--2021) (Table~\ref{tab:cohort}).
The Fudan sub-cohort is a relapse-enriched referral series (all 21
patients relapsed), so the pooled external relapse rate (23 of 41,
56.1\%) sits far above that of an unselected ALL population; external
relapse metrics are therefore exploratory (see Limitations).
Pooling the two hospitals resolves the single-class collapse a
per-site split would otherwise cause for several targets: D19MRD and
Relapse become two-class and AUROC-evaluable across the full cohort,
while D46MRD still has only one positive patient after pooling and is
reported as accuracy.
For external validation, all 301 internal patients were used as the
training set and no external data were used in training.

Feature sets $\mathcal{F}_1$--$\mathcal{F}_7$
(Table~\ref{tab:expdesign}) are temporally ordered and, except for the
$\mathcal{F}_4$ branch, strictly nested: each level adds only features
available at the corresponding clinical timepoint.
$\mathcal{F}_1$ contains WBC count and age (Day~0);
$\mathcal{F}_7$ additionally includes D19MRD and D46MRD values (Day~46).
The feature hierarchy applies identically to all six prediction
targets, with the distinction that classification tasks (targets 1--4)
use AUROC as the primary metric, while survival targets (targets 5--6:
EFS, OS) use Harrell's C-index and time-dependent IPCW AUC.

\begin{table}[ht]
\caption{Cohort characteristics and prediction target distributions.
The External cohort pools patients from two independent hospitals
(see Methods for site details).
$\dagger$~patients with complete labels
for all four classification targets.
$\ddagger$~patients whose automated pipeline output contains at least
one blast cell; only these enter the imaging analysis
(Section~\ref{sec:imgmorph}).}
\label{tab:cohort}
\centering
\small
\begin{tabular}{lcc}
\toprule
\textbf{Characteristic}
  & \textbf{Internal} ($n=301$)
  & \textbf{External} ($n=41$) \\
\midrule
Enrolment years & 2015--2019 & 2017--2023 \\
B-ALL & 163 (54.2\%) & 38 (92.7\%) \\
Patients with complete labels$^\dagger$ & 301 & 29 \\
\midrule
\multicolumn{3}{l}{\textit{Classification targets}} \\
Final Risk: LR & 159 (52.8\%) & 17 (41.5\%) \\
Final Risk: IR/HR & 142 (47.2\%) & 24 (58.5\%) \\
Relapse: Yes & 44 (14.6\%) & 23 (56.1\%) \\
D19MRD $\geq 1\%$ & [X] ([X]\%) & 4 \\
D46MRD $\geq 0.01\%$ & [X] ([X]\%) & 1 \\
\midrule
\multicolumn{3}{l}{\textit{Survival endpoints}} \\
EFS events & 75 (24.9\%) & 24 (58.5\%) \\
OS events & 24 (8.0\%) & 7 (17.1\%) \\
\midrule
\multicolumn{3}{l}{\textit{Multimodal data}} \\
Bone marrow images & [X] (XX\%) & 41 \\
Morphology features & 287 (95.3\%) & 40 \\
\;\,non-empty blast profile$^\ddagger$ & 221 (73.4\%) & 13 \\
\bottomrule
\end{tabular}
\end{table}

%%--------------------------------------------------------------%%
\subsection{Phase~1: LLM zero-shot performance across six targets}
\label{sec:phase1}
%%--------------------------------------------------------------%%

We evaluated four LLMs in a zero-shot setting across all seven
clinical feature levels and all six prediction targets; Table~\ref{tab:phase1_full}
reports the complete internal per-feature-level results and
Table~\ref{tab:phase1_ext} the merged external cohort.
Models included Gemini-3.1-Flash-Lite (Google AI API) and three
open-weight models deployed locally via vLLM: Gemma-4-31B-it,
InternVL3.5-30B-A3B, and Qwen3.6-35B-A3B.
For classification tasks (targets 1--4), each patient's feature vector
was serialised into natural language and the model output a categorical
label; we report AUROC.
For survival targets (5--6), the model output a continuous risk score
(0--1) evaluated by concordance index (C-index).

\begin{table}[!t]
\caption{\textbf{Phase~1 LLM zero-shot performance across all seven clinical feature levels and six prediction targets (internal cohort, $n=301$).} Each cell is the macro-averaged AUROC (classification targets) or Harrell's concordance index (survival targets) for one LLM at one feature level; \textbf{bold} marks each model's best feature level for that target (column maximum). Classification AUROC is computed from parsed categorical predictions (balanced-accuracy form). \textemdash{}~indicates a target not evaluated at that level: D19/D46MRD are used as \emph{input} features at $\mathcal{F}_7$ (see Table~\ref{tab:expdesign}).}
\label{tab:phase1_full}
\centering
\small
\setlength{\tabcolsep}{4pt}
\begin{tabular}{l cccc @{\hskip 1.6em} l cccc}
\toprule
 & \multicolumn{4}{c}{\textbf{Model}} & & \multicolumn{4}{c}{\textbf{Model}} \\
\cmidrule(lr){2-5}\cmidrule(lr){7-10}
\textbf{Lv.} & Gem. & Gma. & Int. & Qw. & \textbf{Lv.} & Gem. & Gma. & Int. & Qw. \\
\midrule
\multicolumn{5}{l}{\textit{Final Risk (AUROC)}} & \multicolumn{5}{l}{\textit{Relapse (AUROC)}} \\
$\mathcal{F}_1$ & 0.788 & 0.785 & 0.729 & 0.792 & $\mathcal{F}_1$ & 0.514 & 0.572 & 0.566 & 0.530 \\
$\mathcal{F}_2$ & 0.782 & 0.796 & 0.539 & 0.789 & $\mathcal{F}_2$ & 0.557 & 0.564 & 0.519 & 0.526 \\
$\mathcal{F}_3$ & 0.812 & 0.842 & 0.583 & 0.831 & $\mathcal{F}_3$ & 0.577 & 0.572 & 0.503 & 0.503 \\
$\mathcal{F}_4$ & 0.826 & 0.835 & 0.591 & 0.837 & $\mathcal{F}_4$ & 0.575 & 0.555 & 0.500 & 0.552 \\
$\mathcal{F}_5$ & 0.791 & 0.854 & 0.823 & 0.845 & $\mathcal{F}_5$ & 0.596 & 0.558 & 0.554 & 0.504 \\
$\mathcal{F}_6$ & 0.834 & \textbf{0.891} & \textbf{0.848} & \textbf{0.882} & $\mathcal{F}_6$ & \textbf{0.599} & 0.552 & \textbf{0.643} & 0.546 \\
$\mathcal{F}_7$ & \textbf{0.849} & 0.788 & 0.716 & 0.856 & $\mathcal{F}_7$ & 0.567 & \textbf{0.586} & 0.540 & \textbf{0.557} \\
\multicolumn{5}{l}{\textit{D19 MRD (AUROC)}} & \multicolumn{5}{l}{\textit{EFS (C-index)}} \\
$\mathcal{F}_1$ & 0.571 & 0.543 & 0.564 & 0.531 & $\mathcal{F}_1$ & 0.577 & 0.591 & 0.584 & 0.573 \\
$\mathcal{F}_2$ & 0.543 & 0.555 & 0.479 & 0.521 & $\mathcal{F}_2$ & 0.607 & 0.598 & 0.522 & 0.561 \\
$\mathcal{F}_3$ & 0.631 & 0.599 & 0.484 & 0.557 & $\mathcal{F}_3$ & 0.587 & 0.614 & 0.545 & 0.537 \\
$\mathcal{F}_4$ & 0.583 & 0.597 & 0.488 & 0.520 & $\mathcal{F}_4$ & 0.598 & 0.614 & 0.526 & 0.561 \\
$\mathcal{F}_5$ & 0.603 & 0.620 & \textbf{0.567} & \textbf{0.605} & $\mathcal{F}_5$ & 0.624 & 0.583 & 0.604 & 0.542 \\
$\mathcal{F}_6$ & \textbf{0.651} & \textbf{0.652} & 0.548 & 0.581 & $\mathcal{F}_6$ & 0.643 & 0.630 & 0.621 & 0.581 \\
$\mathcal{F}_7$ & \textemdash & \textemdash & \textemdash & \textemdash & $\mathcal{F}_7$ & \textbf{0.687} & \textbf{0.716} & \textbf{0.642} & \textbf{0.662} \\
\multicolumn{5}{l}{\textit{D46 MRD (AUROC)}} & \multicolumn{5}{l}{\textit{OS (C-index)}} \\
$\mathcal{F}_1$ & 0.505 & 0.486 & \textbf{0.542} & 0.501 & $\mathcal{F}_1$ & 0.689 & 0.648 & 0.652 & 0.661 \\
$\mathcal{F}_2$ & 0.492 & 0.475 & 0.501 & 0.495 & $\mathcal{F}_2$ & 0.647 & 0.702 & 0.550 & 0.634 \\
$\mathcal{F}_3$ & 0.502 & 0.509 & 0.503 & 0.500 & $\mathcal{F}_3$ & 0.682 & 0.678 & 0.537 & 0.634 \\
$\mathcal{F}_4$ & 0.505 & 0.539 & 0.505 & 0.533 & $\mathcal{F}_4$ & 0.687 & 0.691 & 0.497 & 0.660 \\
$\mathcal{F}_5$ & 0.549 & \textbf{0.577} & 0.518 & \textbf{0.552} & $\mathcal{F}_5$ & 0.701 & 0.704 & 0.734 & 0.619 \\
$\mathcal{F}_6$ & \textbf{0.610} & 0.539 & 0.523 & 0.510 & $\mathcal{F}_6$ & 0.667 & 0.693 & 0.755 & 0.687 \\
$\mathcal{F}_7$ & \textemdash & \textemdash & \textemdash & \textemdash & $\mathcal{F}_7$ & \textbf{0.725} & \textbf{0.718} & \textbf{0.764} & \textbf{0.718} \\
\bottomrule
\end{tabular}
\end{table}

\textbf{Classification performance on the internal cohort.}
Final Risk was the easiest target for all models.
At Day~0 features ($\mathcal{F}_1$: WBC and age), all four models
already achieved AUROC~$\approx 0.73$--$0.79$, reflecting the known
prognostic importance of WBC count.
Performance increased incrementally and peaked at the cytogenetics
panel $\mathcal{F}_6$ (Gemini~0.834, Gemma~0.891, InternVL~0.848,
Qwen~0.882), which was the best level for the three open-weight models;
supplying the D19/D46 MRD values as inputs at $\mathcal{F}_7$ raised
Gemini to 0.849 but did not further improve the others.
Best-level Final Risk AUROCs were thus 0.849 (Gemini,~$\mathcal{F}_7$),
0.891 (Gemma,~$\mathcal{F}_6$), 0.848 (InternVL,~$\mathcal{F}_6$), and
0.882 (Qwen,~$\mathcal{F}_6$).
Relapse prediction was substantially harder: best AUROCs were
0.643 (InternVL,~$\mathcal{F}_6$), 0.599 (Gemini,~$\mathcal{F}_6$),
0.586 (Gemma,~$\mathcal{F}_7$), and 0.557 (Qwen,~$\mathcal{F}_7$),
with no consistent improvement as features accumulated.
D19MRD and D46MRD showed moderate performance (AUROC~0.55--0.65),
suggesting that LLMs can partially infer early treatment response from
immunophenotypic and cytogenetic profiles but that these tasks remain
harder without direct MRD measurement.

\textbf{Survival performance on the internal cohort.}
For EFS, best-level C-indices were all attained at $\mathcal{F}_7$:
0.716 (Gemma), 0.687 (Gemini), 0.662 (Qwen), and 0.642 (InternVL).
The inclusion of the D19/D46 MRD values at $\mathcal{F}_7$ was
consistently optimal for EFS, confirming that residual disease burden
encodes strong prognostic information accessible to LLMs without
labelled training data.
For OS—with only 24 events—C-indices also peaked at $\mathcal{F}_7$:
InternVL reached 0.764, Gemini 0.725, and Gemma and Qwen both 0.718.
Log-rank tests were significant ($p < 0.05$) at $\mathcal{F}_7$
for EFS for all four LLMs.

\textbf{External validation (merged cohort).}
Final Risk transferred well on the pooled 41-patient external cohort:
LLM zero-shot AUROC ranged 0.450--0.925 across models and feature
levels, peaking at $\mathcal{F}_7$ (0.750--0.925 across the four
LLMs), close to the internal range and obtained without any
cohort-specific retraining (Table~\ref{tab:phase1_ext}).
Pooling the two source hospitals resolves the single-class collapse
that a per-site split would otherwise produce for Relapse and D19MRD:
both targets are now two-class across the full external cohort and
AUROC-evaluable (Relapse 0.482--0.717; D19MRD close to chance,
0.340--0.550); only D46MRD still has a single positive patient after
pooling (1 of 29 labelled patients) and is reported with accuracy
(0.035--0.724) rather than AUROC.
For survival, the pooled cohort (24 EFS events, 7 deaths) gave LLM EFS
C-indices of 0.542--0.715 and OS 0.548--0.859.

\begin{table}[p]
\caption{\textbf{Phase~1 LLM zero-shot on the pooled external cohort,
all feature levels} (Gm/Gma/Iv/Qw = Gemini/Gemma/InternVL/Qwen).
The two source hospitals are pooled into a
single external cohort ($n$=41 for survival; classification targets
use the $n$=29 subset with available labels). Classification: AUROC
($^*$~= accuracy; D46MRD has only one positive class after pooling,
so accuracy is reported instead). Survival (EFS, OS): C-index.
\textbf{Bold}: best feature level per column within each task block
(ties all bolded). \textemdash: not evaluable ($\mathcal{F}_7$ not
defined for D19/D46 MRD). External pooled cohort; exploratory,
indicative only.}
\label{tab:phase1_ext}
\centering
\footnotesize
\renewcommand{\arraystretch}{0.8}
\setlength{\tabcolsep}{5pt}
\begin{tabular}{@{}l cccc@{}}
\toprule
\textbf{Lv.} & Gm & Gma & Iv & Qw \\
\midrule
\multicolumn{5}{l}{\textit{Final Risk (AUROC)}} \\
$\mathcal{F}_1$ & 0.739 & 0.794 & 0.714 & 0.794 \\
$\mathcal{F}_2$ & 0.794 & 0.819 & 0.450 & 0.794 \\
$\mathcal{F}_3$ & 0.708 & 0.819 & 0.531 & 0.794 \\
$\mathcal{F}_4$ & 0.844 & 0.819 & 0.697 & 0.794 \\
$\mathcal{F}_5$ & 0.653 & \textbf{0.894} & 0.536 & 0.739 \\
$\mathcal{F}_6$ & 0.708 & 0.844 & 0.733 & 0.819 \\
$\mathcal{F}_7$ & \textbf{0.925} & 0.775 & \textbf{0.750} & \textbf{0.850} \\
\multicolumn{5}{l}{\textit{D19 MRD (AUROC)}} \\
$\mathcal{F}_1$ & 0.425 & 0.425 & 0.490 & 0.380 \\
$\mathcal{F}_2$ & 0.530 & 0.405 & 0.500 & \textbf{0.550} \\
$\mathcal{F}_3$ & 0.425 & 0.405 & 0.500 & 0.340 \\
$\mathcal{F}_4$ & \textbf{0.550} & 0.405 & 0.500 & 0.445 \\
$\mathcal{F}_5$ & 0.425 & \textbf{0.465} & 0.395 & 0.445 \\
$\mathcal{F}_6$ & 0.425 & 0.405 & \textbf{0.515} & 0.545 \\
$\mathcal{F}_7$ & \textemdash & \textemdash & \textemdash & \textemdash \\
\multicolumn{5}{l}{\textit{D46 MRD (accuracy; single positive)}} \\
$\mathcal{F}_1$ & \textbf{0.690$^*$} & 0.586$^*$ & \textbf{0.586$^*$} & 0.621$^*$ \\
$\mathcal{F}_2$ & 0.586$^*$ & 0.552$^*$ & 0.035$^*$ & 0.655$^*$ \\
$\mathcal{F}_3$ & 0.586$^*$ & 0.586$^*$ & 0.035$^*$ & 0.690$^*$ \\
$\mathcal{F}_4$ & 0.552$^*$ & 0.586$^*$ & 0.035$^*$ & \textbf{0.724$^*$} \\
$\mathcal{F}_5$ & 0.586$^*$ & \textbf{0.621$^*$} & 0.138$^*$ & 0.690$^*$ \\
$\mathcal{F}_6$ & 0.621$^*$ & 0.586$^*$ & 0.207$^*$ & 0.655$^*$ \\
$\mathcal{F}_7$ & \textemdash & \textemdash & \textemdash & \textemdash \\
\multicolumn{5}{l}{\textit{Relapse (AUROC)}} \\
$\mathcal{F}_1$ & 0.587 & 0.634 & \textbf{0.699} & 0.482 \\
$\mathcal{F}_2$ & 0.504 & 0.612 & 0.500 & 0.482 \\
$\mathcal{F}_3$ & 0.674 & 0.551 & 0.562 & 0.544 \\
$\mathcal{F}_4$ & 0.609 & 0.551 & 0.685 & 0.565 \\
$\mathcal{F}_5$ & \textbf{0.717} & 0.612 & 0.620 & 0.587 \\
$\mathcal{F}_6$ & 0.634 & \textbf{0.656} & 0.536 & \textbf{0.609} \\
$\mathcal{F}_7$ & 0.500 & 0.522 & 0.500 & 0.522 \\
\multicolumn{5}{l}{\textit{EFS (C-index)}} \\
$\mathcal{F}_1$ & 0.657 & \textbf{0.669} & 0.627 & 0.634 \\
$\mathcal{F}_2$ & \textbf{0.715} & 0.659 & 0.551 & \textbf{0.645} \\
$\mathcal{F}_3$ & 0.635 & 0.590 & 0.615 & 0.638 \\
$\mathcal{F}_4$ & 0.634 & 0.590 & 0.615 & 0.638 \\
$\mathcal{F}_5$ & 0.617 & 0.656 & 0.639 & 0.542 \\
$\mathcal{F}_6$ & 0.657 & 0.651 & \textbf{0.709} & 0.567 \\
$\mathcal{F}_7$ & 0.656 & 0.664 & 0.702 & 0.591 \\
\multicolumn{5}{l}{\textit{OS (C-index)}} \\
$\mathcal{F}_1$ & 0.808 & \textbf{0.821} & 0.802 & 0.714 \\
$\mathcal{F}_2$ & 0.804 & 0.792 & 0.633 & 0.744 \\
$\mathcal{F}_3$ & 0.746 & 0.780 & 0.710 & 0.548 \\
$\mathcal{F}_4$ & 0.760 & 0.780 & 0.710 & 0.548 \\
$\mathcal{F}_5$ & 0.762 & 0.808 & 0.816 & 0.694 \\
$\mathcal{F}_6$ & \textbf{0.810} & 0.603 & 0.814 & \textbf{0.806} \\
$\mathcal{F}_7$ & 0.802 & 0.615 & \textbf{0.859} & 0.665 \\
\bottomrule
\end{tabular}
\end{table}

%%--------------------------------------------------------------%%
\subsection{Testing whether bone marrow imaging adds prognostic
information beyond clinical features}
\label{sec:imgmorph}
%%--------------------------------------------------------------%%

Bone marrow smear microscopy is performed for every ALL patient at
diagnosis, yet current risk-stratification rules use only the blast
percentage derived from it—the images themselves enter no downstream
risk calculation, and prior computational studies of paediatric ALL
risk have likewise not tested whether the images carry additional
prognostic signal.
We therefore evaluate, at every clinical feature level
$\mathcal{F}_3$--$\mathcal{F}_7$, three ways of injecting imaging
information into the zero-shot LLMs:
($+$img)~the raw smear images as direct multimodal input;
($+$morph)~a structured morphology profile generated by a four-stage
automated pipeline (SAM segmentation, $\sim$37 morphological features
per cell, BiomedCLIP cell-type clustering, per-patient aggregation;
Methods, Section~\ref{sec:morph}), serialised as text;
($+$both)~images and profile together.
To keep comparisons fair, each level's base row is re-evaluated on the
same patient subset as its imaging variants (patients with available
morphology data).
Table~\ref{tab:imgmorph_int} reports the internal cohort;
Tables~\ref{tab:imgmorph_ext_cls} and~\ref{tab:imgmorph_ext_surv}
report the merged external cohort.

With the comparison now complete at all five levels, the clearest and
most consistent internal-cohort gain is for relapse prediction from
the morphology profile: $+$morph improves Relapse AUROC over base for
a majority of models (3--4 of 4) at three of the five levels (4/4 at
$\mathcal{F}_3$ and $\mathcal{F}_5$, 3/4 at $\mathcal{F}_6$), with
per-model deltas ranging from $-0.043$ to $+0.070$ across levels.
Raw images help less reliably: $+$img improves Final Risk for 3/4
models at $\mathcal{F}_3$ and $\mathcal{F}_6$ but for at most 2/4 at
the other levels, and improves Relapse for only 1/4 models at
$\mathcal{F}_3$--$\mathcal{F}_5$ (4/4 at $\mathcal{F}_6$).
At $\mathcal{F}_7$, where the D19/D46 MRD measurements are already
inputs, the morphology profile still adds Final Risk signal—improving
all four models ($+0.005$ to $+0.102$)—whereas raw images degrade
Final Risk for all four ($-0.034$ to $-0.005$); Relapse changes at
$\mathcal{F}_7$ are small for every variant ($-0.019$ to $+0.027$).
For survival, $+$morph is again the only variant with a consistent
benefit: it improves EFS over base in 15 of 20 model--level
comparisons (4/4 models at $\mathcal{F}_4$, 3/4 at
$\mathcal{F}_5$--$\mathcal{F}_7$; deltas $-0.027$ to $+0.033$),
whereas $+$img (9/20) and $+$both (6/20) are net neutral to harmful,
and for OS no variant helps consistently (at best $+$morph, 8/20;
$+$img degrades OS by up to $-0.235$).
The morphology profile also augments the tabular Phase~2 random
forest, but systematically across all five clinical levels
($\mathcal{F}_3$--$\mathcal{F}_7$, Table~\ref{tab:phase2_morph}) the
effect is not uniformly positive. On the internal temporal splits,
adding the SHAP-selected morphology features leaves classification
essentially unchanged on average (mean $\Delta$AUROC~$\approx-0.003$,
Wilcoxon $p=0.86$) and significantly worsens survival discrimination
(mean $\Delta$C-index~$=-0.052$, $p<0.001$; morph better in only
8\% of level--task--split cells), consistent with the extra features
overfitting the small temporal splits. On the pooled external
cohort the pattern is different for survival: RSF C-index improves
with morphology in 80\% of cells (mean $\Delta=+0.087$, $p=0.027$),
while external classification stays neutral (mean
$\Delta\approx-0.002$) — echoing the Phase~1 finding that morphology
aids survival more than classification.
The $+$img/$+$morph/$+$both comparisons above are unadjusted differences on a single internal cohort, without
significance testing, and should be read as directional evidence; the
external fair-comparison subsets (Tables~\ref{tab:imgmorph_ext_cls}
and~\ref{tab:imgmorph_ext_surv}) are too small to confirm or refute
them.

\begin{table}[p]
\caption{\textbf{Information gain from bone marrow imaging at every
clinical feature level (internal cohort).} For each base level
$\mathcal{F}_k$ ($k=3..7$): $+$img adds raw smear images, $+$morph the
structured morphology profile, $+$both adds both. Classification
targets: AUROC; EFS/OS: C-index. Base rows are re-evaluated on the
subset of patients with a non-empty automated morphology profile
(blast cells detected by the pipeline), so they can differ slightly
from Table~\ref{tab:phase1_full}. \textbf{Bold}: best variant within
each level block (per model column; ties all bolded). \textemdash: not
applicable (D19/D46MRD are inputs at $\mathcal{F}_7$).}
\label{tab:imgmorph_int}
\centering
\tiny
\renewcommand{\arraystretch}{0.5}
\setlength{\tabcolsep}{3.5pt}
\begin{tabular}{@{}l cccc @{\hskip 1.2em} l cccc@{}}
\toprule
 & Gm & Gma & Iv & Qw & & Gm & Gma & Iv & Qw \\
\midrule
\multicolumn{5}{l}{\textit{Final Risk (AUROC)}} & \multicolumn{5}{l}{\textit{Relapse (AUROC)}} \\
$\mathcal{F}_3$      & 0.792 & 0.827 & 0.558 & 0.812 & $\mathcal{F}_3$      & 0.543 & 0.532 & 0.505 & 0.481 \\
\;\,$+$img           & \textbf{0.801} & 0.823 & \textbf{0.561} & \textbf{0.822} & \;\,$+$img           & 0.484 & 0.513 & 0.484 & 0.486 \\
\;\,$+$morph         & 0.759 & \textbf{0.833} & 0.546 & 0.798 & \;\,$+$morph         & \textbf{0.565} & \textbf{0.535} & \textbf{0.511} & \textbf{0.521} \\
\;\,$+$both          & 0.767 & 0.811 & 0.552 & 0.769 & \;\,$+$both          & 0.554 & 0.497 & 0.503 & 0.510 \\
\addlinespace[2pt]
$\mathcal{F}_4$      & \textbf{0.811} & 0.822 & 0.560 & 0.817 & $\mathcal{F}_4$      & 0.540 & 0.511 & 0.484 & \textbf{0.536} \\
\;\,$+$img           & 0.757 & \textbf{0.823} & 0.579 & 0.809 & \;\,$+$img           & 0.508 & \textbf{0.524} & 0.465 & 0.490 \\
\;\,$+$morph         & 0.782 & 0.821 & \textbf{0.616} & \textbf{0.839} & \;\,$+$morph         & 0.557 & 0.468 & 0.511 & 0.515 \\
\;\,$+$both          & 0.776 & 0.794 & 0.577 & 0.800 & \;\,$+$both          & \textbf{0.562} & 0.508 & \textbf{0.516} & 0.508 \\
\addlinespace[2pt]
$\mathcal{F}_5$      & 0.783 & 0.839 & 0.804 & 0.827 & $\mathcal{F}_5$      & 0.602 & 0.546 & 0.540 & 0.502 \\
\;\,$+$img           & 0.763 & 0.806 & \textbf{0.840} & 0.812 & \;\,$+$img           & 0.586 & 0.513 & 0.557 & 0.491 \\
\;\,$+$morph         & 0.784 & 0.779 & 0.795 & \textbf{0.843} & \;\,$+$morph         & \textbf{0.616} & \textbf{0.594} & \textbf{0.581} & 0.550 \\
\;\,$+$both          & \textbf{0.788} & \textbf{0.844} & 0.775 & 0.795 & \;\,$+$both          & \textbf{0.616} & 0.551 & 0.530 & \textbf{0.572} \\
\addlinespace[2pt]
$\mathcal{F}_6$      & 0.818 & 0.872 & 0.832 & 0.874 & $\mathcal{F}_6$      & 0.546 & 0.479 & 0.618 & 0.486 \\
\;\,$+$img           & \textbf{0.840} & 0.880 & \textbf{0.838} & 0.873 & \;\,$+$img           & \textbf{0.624} & 0.495 & \textbf{0.621} & \textbf{0.505} \\
\;\,$+$morph         & 0.820 & 0.875 & 0.800 & \textbf{0.876} & \;\,$+$morph         & 0.616 & \textbf{0.527} & 0.602 & 0.492 \\
\;\,$+$both          & 0.821 & \textbf{0.881} & 0.822 & 0.862 & \;\,$+$both          & 0.600 & 0.508 & 0.565 & 0.502 \\
\addlinespace[2pt]
$\mathcal{F}_7$      & 0.837 & 0.761 & 0.699 & 0.852 & $\mathcal{F}_7$      & \textbf{0.589} & 0.591 & 0.559 & 0.564 \\
\;\,$+$img           & 0.803 & 0.745 & 0.694 & 0.832 & \;\,$+$img           & \textbf{0.589} & \textbf{0.597} & 0.567 & \textbf{0.578} \\
\;\,$+$morph         & \textbf{0.842} & \textbf{0.770} & \textbf{0.801} & \textbf{0.858} & \;\,$+$morph         & \textbf{0.589} & 0.573 & 0.557 & 0.564 \\
\;\,$+$both          & 0.786 & 0.731 & 0.715 & 0.765 & \;\,$+$both          & 0.575 & 0.581 & \textbf{0.586} & \textbf{0.578} \\
\addlinespace[2pt]
Avg & \textbf{0.808} & \textbf{0.824} & 0.691 & 0.836 & Avg & 0.564 & 0.532 & 0.541 & 0.514 \\
\;\,$+$img & 0.793 & 0.815 & 0.702 & 0.830 & \;\,$+$img & 0.558 & 0.528 & 0.539 & 0.510 \\
\;\,$+$morph & 0.797 & 0.816 & \textbf{0.712} & \textbf{0.843} & \;\,$+$morph & \textbf{0.588} & \textbf{0.539} & \textbf{0.552} & 0.529 \\
\;\,$+$both & 0.787 & 0.812 & 0.688 & 0.798 & \;\,$+$both & 0.581 & 0.529 & 0.540 & \textbf{0.534} \\
\midrule
\multicolumn{5}{l}{\textit{D19 MRD (AUROC)}} & \multicolumn{5}{l}{\textit{EFS (C-index)}} \\
$\mathcal{F}_3$      & \textbf{0.615} & 0.608 & 0.478 & 0.575 & $\mathcal{F}_3$      & 0.550 & 0.585 & \textbf{0.513} & 0.532 \\
\;\,$+$img           & 0.590 & \textbf{0.704} & 0.499 & \textbf{0.644} & \;\,$+$img           & 0.580 & 0.575 & 0.448 & \textbf{0.552} \\
\;\,$+$morph         & 0.538 & 0.677 & 0.488 & 0.412 & \;\,$+$morph         & 0.567 & \textbf{0.593} & 0.496 & 0.525 \\
\;\,$+$both          & 0.579 & 0.643 & \textbf{0.512} & 0.461 & \;\,$+$both          & \textbf{0.602} & 0.576 & 0.452 & 0.522 \\
\addlinespace[2pt]
$\mathcal{F}_4$      & 0.591 & 0.606 & 0.478 & 0.519 & $\mathcal{F}_4$      & 0.559 & 0.584 & 0.484 & 0.555 \\
\;\,$+$img           & \textbf{0.626} & 0.642 & \textbf{0.508} & 0.548 & \;\,$+$img           & \textbf{0.612} & \textbf{0.600} & 0.456 & \textbf{0.572} \\
\;\,$+$morph         & 0.609 & \textbf{0.658} & 0.493 & \textbf{0.613} & \;\,$+$morph         & 0.568 & 0.585 & \textbf{0.502} & 0.560 \\
\;\,$+$both          & 0.584 & 0.586 & \textbf{0.508} & 0.536 & \;\,$+$both          & 0.577 & 0.589 & 0.470 & 0.547 \\
\addlinespace[2pt]
$\mathcal{F}_5$      & 0.612 & 0.632 & \textbf{0.574} & \textbf{0.644} & $\mathcal{F}_5$      & 0.608 & 0.552 & \textbf{0.574} & 0.520 \\
\;\,$+$img           & 0.615 & 0.636 & 0.504 & 0.606 & \;\,$+$img           & 0.585 & 0.577 & 0.473 & 0.500 \\
\;\,$+$morph         & \textbf{0.655} & 0.608 & 0.538 & 0.611 & \;\,$+$morph         & \textbf{0.609} & \textbf{0.585} & 0.574 & \textbf{0.535} \\
\;\,$+$both          & 0.644 & \textbf{0.655} & 0.501 & 0.589 & \;\,$+$both          & 0.584 & 0.582 & 0.570 & 0.503 \\
\addlinespace[2pt]
$\mathcal{F}_6$      & 0.678 & 0.669 & 0.560 & 0.606 & $\mathcal{F}_6$      & 0.625 & \textbf{0.613} & 0.606 & 0.547 \\
\;\,$+$img           & 0.660 & \textbf{0.694} & \textbf{0.590} & \textbf{0.679} & \;\,$+$img           & 0.602 & 0.598 & 0.556 & 0.550 \\
\;\,$+$morph         & 0.636 & 0.684 & 0.556 & 0.591 & \;\,$+$morph         & \textbf{0.650} & 0.586 & \textbf{0.610} & 0.577 \\
\;\,$+$both          & \textbf{0.720} & 0.684 & 0.563 & 0.627 & \;\,$+$both          & 0.624 & 0.585 & 0.578 & \textbf{0.593} \\
\addlinespace[2pt]
$\mathcal{F}_7$      & \textemdash & \textemdash & \textemdash & \textemdash & $\mathcal{F}_7$      & 0.658 & 0.700 & 0.615 & 0.622 \\
\;\,$+$img           & \textemdash & \textemdash & \textemdash & \textemdash & \;\,$+$img           & 0.659 & 0.668 & 0.518 & \textbf{0.631} \\
\;\,$+$morph         & \textemdash & \textemdash & \textemdash & \textemdash & \;\,$+$morph         & \textbf{0.667} & \textbf{0.701} & \textbf{0.636} & 0.621 \\
\;\,$+$both          & \textemdash & \textemdash & \textemdash & \textemdash & \;\,$+$both          & 0.666 & 0.650 & 0.545 & 0.599 \\
\addlinespace[2pt]
Avg & 0.624 & 0.629 & 0.522 & 0.586 & Avg & 0.600 & 0.607 & 0.558 & 0.555 \\
\;\,$+$img & 0.623 & \textbf{0.669} & \textbf{0.525} & \textbf{0.619} & \;\,$+$img & 0.608 & 0.604 & 0.490 & 0.561 \\
\;\,$+$morph & 0.609 & 0.657 & 0.519 & 0.557 & \;\,$+$morph & \textbf{0.612} & \textbf{0.610} & \textbf{0.563} & \textbf{0.564} \\
\;\,$+$both & \textbf{0.632} & 0.642 & 0.521 & 0.553 & \;\,$+$both & 0.610 & 0.596 & 0.523 & 0.553 \\
\midrule
\multicolumn{5}{l}{\textit{D46 MRD (AUROC)}} & \multicolumn{5}{l}{\textit{OS (C-index)}} \\
$\mathcal{F}_3$      & 0.501 & 0.488 & 0.498 & 0.488 & $\mathcal{F}_3$      & 0.684 & 0.630 & \textbf{0.482} & 0.622 \\
\;\,$+$img           & 0.516 & 0.524 & \textbf{0.506} & 0.498 & \;\,$+$img           & 0.672 & 0.639 & 0.343 & 0.584 \\
\;\,$+$morph         & 0.535 & 0.505 & 0.505 & 0.519 & \;\,$+$morph         & 0.634 & 0.633 & 0.460 & 0.595 \\
\;\,$+$both          & \textbf{0.560} & \textbf{0.589} & 0.420 & \textbf{0.547} & \;\,$+$both          & \textbf{0.687} & \textbf{0.647} & 0.430 & \textbf{0.637} \\
\addlinespace[2pt]
$\mathcal{F}_4$      & \textbf{0.509} & 0.529 & 0.495 & \textbf{0.558} & $\mathcal{F}_4$      & 0.705 & \textbf{0.663} & \textbf{0.439} & 0.639 \\
\;\,$+$img           & 0.501 & 0.551 & 0.508 & 0.546 & \;\,$+$img           & 0.672 & 0.635 & 0.387 & 0.576 \\
\;\,$+$morph         & 0.505 & 0.546 & 0.477 & 0.485 & \;\,$+$morph         & 0.660 & 0.652 & 0.430 & \textbf{0.677} \\
\;\,$+$both          & 0.502 & \textbf{0.582} & \textbf{0.521} & 0.493 & \;\,$+$both          & \textbf{0.714} & 0.649 & 0.366 & 0.600 \\
\addlinespace[2pt]
$\mathcal{F}_5$      & 0.525 & 0.565 & 0.494 & 0.560 & $\mathcal{F}_5$      & 0.695 & \textbf{0.701} & \textbf{0.720} & 0.602 \\
\;\,$+$img           & \textbf{0.626} & 0.504 & 0.529 & 0.496 & \;\,$+$img           & 0.682 & 0.622 & 0.485 & 0.509 \\
\;\,$+$morph         & 0.598 & 0.508 & \textbf{0.542} & 0.488 & \;\,$+$morph         & \textbf{0.728} & 0.670 & 0.712 & \textbf{0.620} \\
\;\,$+$both          & 0.578 & \textbf{0.570} & 0.466 & \textbf{0.572} & \;\,$+$both          & 0.683 & 0.602 & 0.691 & 0.530 \\
\addlinespace[2pt]
$\mathcal{F}_6$      & 0.606 & 0.524 & \textbf{0.536} & 0.503 & $\mathcal{F}_6$      & 0.658 & 0.683 & \textbf{0.742} & \textbf{0.665} \\
\;\,$+$img           & 0.604 & 0.565 & 0.499 & 0.497 & \;\,$+$img           & 0.656 & 0.633 & 0.580 & 0.588 \\
\;\,$+$morph         & \textbf{0.685} & \textbf{0.592} & 0.487 & \textbf{0.528} & \;\,$+$morph         & 0.675 & \textbf{0.701} & 0.704 & 0.638 \\
\;\,$+$both          & 0.679 & 0.531 & 0.511 & 0.471 & \;\,$+$both          & \textbf{0.680} & 0.658 & 0.680 & 0.583 \\
\addlinespace[2pt]
$\mathcal{F}_7$      & \textemdash & \textemdash & \textemdash & \textemdash & $\mathcal{F}_7$      & 0.680 & 0.698 & \textbf{0.752} & 0.668 \\
\;\,$+$img           & \textemdash & \textemdash & \textemdash & \textemdash & \;\,$+$img           & 0.669 & \textbf{0.718} & 0.532 & 0.677 \\
\;\,$+$morph         & \textemdash & \textemdash & \textemdash & \textemdash & \;\,$+$morph         & \textbf{0.702} & 0.693 & 0.739 & \textbf{0.685} \\
\;\,$+$both          & \textemdash & \textemdash & \textemdash & \textemdash & \;\,$+$both          & 0.699 & 0.694 & 0.639 & 0.678 \\
\addlinespace[2pt]
Avg & 0.535 & 0.526 & 0.506 & \textbf{0.527} & Avg & 0.684 & \textbf{0.675} & \textbf{0.627} & 0.639 \\
\;\,$+$img & 0.561 & 0.536 & \textbf{0.510} & 0.509 & \;\,$+$img & 0.670 & 0.649 & 0.465 & 0.587 \\
\;\,$+$morph & \textbf{0.581} & 0.538 & 0.503 & 0.505 & \;\,$+$morph & 0.680 & 0.670 & 0.609 & \textbf{0.643} \\
\;\,$+$both & 0.580 & \textbf{0.568} & 0.479 & 0.521 & \;\,$+$both & \textbf{0.692} & 0.650 & 0.561 & 0.605 \\
\bottomrule
\end{tabular}
\end{table}

\begin{table}[p]
\caption{\textbf{Information gain from bone marrow imaging at every
clinical feature level, external classification tasks (merged external
cohort).} Gm/Gma/Iv/Qw =
Gemini/Gemma/InternVL/Qwen. For each base level $\mathcal{F}_k$
($k=3..7$ for Final Risk and Relapse, $k=3..6$ for D19/D46MRD): $+$img
adds raw smear images, $+$morph the structured morphology profile,
$+$both adds both. AUROC ($^*$~= accuracy, single-class task/level).
Base rows are re-evaluated on the morphology-available patient subset,
so they can differ slightly from Table~\ref{tab:phase1_ext}.
\textbf{Bold}: best variant within each level block (per column; ties
all bolded). The pooled external fair-comparison subset is small
($n{=}8$; one model--level cell $n{=}7$ owing to a parse-filtered
response); indicative only.}
\label{tab:imgmorph_ext_cls}
\centering
\tiny
\renewcommand{\arraystretch}{0.5}
\setlength{\tabcolsep}{28pt}
\setlength{\aboverulesep}{0.5pt}
\setlength{\belowrulesep}{0.5pt}
\begin{tabular}{@{}l cccc@{}}
\toprule
 & Gm & Gma & Iv & Qw \\
\midrule
\multicolumn{5}{l}{\textit{Final Risk (AUROC; $^*$=accuracy)}} \\
$\mathcal{F}_3$      & \textbf{0.750} & 0.750 & \textbf{0.625} & 0.750 \\
\;\,$+$img           & \textbf{0.750} & 0.750 & 0.500 & 0.750 \\
\;\,$+$morph         & 0.625 & 0.750 & \textbf{0.625} & \textbf{0.875} \\
\;\,$+$both          & \textbf{0.750} & \textbf{0.875} & \textbf{0.625} & 0.500 \\
\addlinespace[0.3pt]
$\mathcal{F}_4$      & \textbf{0.750} & \textbf{0.750} & \textbf{0.750} & \textbf{0.750} \\
\;\,$+$img           & \textbf{0.750} & \textbf{0.750} & 0.625 & \textbf{0.750} \\
\;\,$+$morph         & \textbf{0.750} & \textbf{0.750} & \textbf{0.750} & \textbf{0.750} \\
\;\,$+$both          & 0.625 & \textbf{0.750} & 0.625 & 0.625 \\
\addlinespace[0.3pt]
$\mathcal{F}_5$      & \textbf{0.875} & \textbf{0.875} & 0.625 & 0.750 \\
\;\,$+$img           & 0.750 & \textbf{0.875} & \textbf{0.750} & 0.750 \\
\;\,$+$morph         & 0.750 & 0.750 & \textbf{0.750} & \textbf{0.875} \\
\;\,$+$both          & \textbf{0.875} & \textbf{0.875} & \textbf{0.750} & \textbf{0.875} \\
\addlinespace[0.3pt]
$\mathcal{F}_6$      & \textbf{0.875} & \textbf{0.875} & \textbf{0.875} & \textbf{0.875} \\
\;\,$+$img           & \textbf{0.875} & \textbf{0.875} & \textbf{0.875} & \textbf{0.875} \\
\;\,$+$morph         & \textbf{0.875} & \textbf{0.875} & \textbf{0.875} & \textbf{0.875} \\
\;\,$+$both          & \textbf{0.875} & \textbf{0.875} & \textbf{0.875} & \textbf{0.875} \\
\addlinespace[0.3pt]
$\mathcal{F}_7$      & \textbf{1.000} & 0.750 & 0.750 & \textbf{0.875} \\
\;\,$+$img           & \textbf{1.000} & \textbf{0.875} & \textbf{0.875} & 0.750 \\
\;\,$+$morph         & \textbf{1.000} & 0.625 & 0.750 & \textbf{0.875} \\
\;\,$+$both          & 0.750 & 0.750 & 0.625 & \textbf{0.875} \\
\addlinespace[2pt]
Avg & \textbf{0.850} & 0.800 & 0.725 & 0.800 \\
\;\,$+$img & 0.825 & \textbf{0.825} & 0.725 & 0.775 \\
\;\,$+$morph & 0.800 & 0.750 & \textbf{0.750} & \textbf{0.850} \\
\;\,$+$both & 0.775 & \textbf{0.825} & 0.700 & 0.750 \\
\midrule
\multicolumn{5}{l}{\textit{D19 MRD (AUROC; $^*$=accuracy)}} \\
$\mathcal{F}_3$      & \textbf{0.625$^*$} & \textbf{0.625$^*$} & 0.000$^*$ & \textbf{0.625$^*$} \\
\;\,$+$img           & \textbf{0.625$^*$} & 0.500$^*$ & 0.000$^*$ & 0.500$^*$ \\
\;\,$+$morph         & 0.375$^*$ & 0.125$^*$ & 0.000$^*$ & \textbf{0.625$^*$} \\
\;\,$+$both          & 0.375$^*$ & 0.000$^*$ & \textbf{0.125$^*$} & 0.375$^*$ \\
\addlinespace[0.3pt]
$\mathcal{F}_4$      & \textbf{0.625$^*$} & \textbf{0.625$^*$} & \textbf{0.000$^*$} & \textbf{0.625$^*$} \\
\;\,$+$img           & \textbf{0.625$^*$} & 0.500$^*$ & \textbf{0.000$^*$} & 0.250$^*$ \\
\;\,$+$morph         & 0.375$^*$ & 0.125$^*$ & \textbf{0.000$^*$} & 0.500$^*$ \\
\;\,$+$both          & 0.375$^*$ & 0.125$^*$ & \textbf{0.000$^*$} & 0.375$^*$ \\
\addlinespace[0.3pt]
$\mathcal{F}_5$      & \textbf{0.875$^*$} & \textbf{0.750$^*$} & 0.000$^*$ & \textbf{0.750$^*$} \\
\;\,$+$img           & \textbf{0.875$^*$} & \textbf{0.750$^*$} & 0.000$^*$ & \textbf{0.750$^*$} \\
\;\,$+$morph         & \textbf{0.875$^*$} & 0.375$^*$ & \textbf{0.250$^*$} & \textbf{0.750$^*$} \\
\;\,$+$both          & 0.750$^*$ & 0.375$^*$ & 0.000$^*$ & \textbf{0.750$^*$} \\
\addlinespace[0.3pt]
$\mathcal{F}_6$      & \textbf{0.750$^*$} & \textbf{0.500$^*$} & 0.375$^*$ & \textbf{1.000$^*$} \\
\;\,$+$img           & \textbf{0.750$^*$} & \textbf{0.500$^*$} & \textbf{0.500$^*$} & 0.500$^*$ \\
\;\,$+$morph         & 0.625$^*$ & 0.375$^*$ & 0.250$^*$ & 0.750$^*$ \\
\;\,$+$both          & 0.625$^*$ & \textbf{0.500$^*$} & 0.250$^*$ & 0.875$^*$ \\
\addlinespace[2pt]
Avg & \textbf{0.719$^*$} & \textbf{0.625$^*$} & 0.094$^*$ & \textbf{0.750$^*$} \\
\;\,$+$img & \textbf{0.719$^*$} & 0.563$^*$ & \textbf{0.125$^*$} & 0.500$^*$ \\
\;\,$+$morph & 0.563$^*$ & 0.250$^*$ & \textbf{0.125$^*$} & 0.656$^*$ \\
\;\,$+$both & 0.531$^*$ & 0.250$^*$ & 0.094$^*$ & 0.594$^*$ \\
\midrule
\multicolumn{5}{l}{\textit{D46 MRD (accuracy; single positive)}} \\
$\mathcal{F}_3$      & \textbf{0.625$^*$} & \textbf{0.625$^*$} & 0.000$^*$ & \textbf{0.625$^*$} \\
\;\,$+$img           & \textbf{0.625$^*$} & 0.500$^*$ & \textbf{0.125$^*$} & 0.500$^*$ \\
\;\,$+$morph         & 0.500$^*$ & 0.125$^*$ & 0.000$^*$ & 0.250$^*$ \\
\;\,$+$both          & 0.500$^*$ & 0.125$^*$ & \textbf{0.125$^*$} & 0.375$^*$ \\
\addlinespace[0.3pt]
$\mathcal{F}_4$      & \textbf{0.625$^*$} & \textbf{0.625$^*$} & \textbf{0.000$^*$} & \textbf{0.750$^*$} \\
\;\,$+$img           & 0.500$^*$ & \textbf{0.625$^*$} & \textbf{0.000$^*$} & 0.500$^*$ \\
\;\,$+$morph         & 0.375$^*$ & 0.250$^*$ & \textbf{0.000$^*$} & 0.625$^*$ \\
\;\,$+$both          & 0.500$^*$ & 0.000$^*$ & \textbf{0.000$^*$} & 0.500$^*$ \\
\addlinespace[0.3pt]
$\mathcal{F}_5$      & \textbf{0.875$^*$} & \textbf{0.625$^*$} & \textbf{0.125$^*$} & \textbf{0.875$^*$} \\
\;\,$+$img           & 0.750$^*$ & \textbf{0.625$^*$} & \textbf{0.125$^*$} & 0.750$^*$ \\
\;\,$+$morph         & \textbf{0.875$^*$} & 0.375$^*$ & 0.000$^*$ & 0.750$^*$ \\
\;\,$+$both          & \textbf{0.875$^*$} & 0.125$^*$ & 0.000$^*$ & 0.750$^*$ \\
\addlinespace[0.3pt]
$\mathcal{F}_6$      & \textbf{0.750$^*$} & \textbf{0.500$^*$} & 0.250$^*$ & \textbf{0.875$^*$} \\
\;\,$+$img           & \textbf{0.750$^*$} & \textbf{0.500$^*$} & \textbf{0.375$^*$} & 0.750$^*$ \\
\;\,$+$morph         & 0.500$^*$ & \textbf{0.500$^*$} & 0.125$^*$ & 0.750$^*$ \\
\;\,$+$both          & 0.625$^*$ & 0.375$^*$ & 0.000$^*$ & 0.625$^*$ \\
\addlinespace[2pt]
Avg & \textbf{0.719$^*$} & \textbf{0.594$^*$} & 0.094$^*$ & \textbf{0.781$^*$} \\
\;\,$+$img & 0.656$^*$ & 0.563$^*$ & \textbf{0.156$^*$} & 0.625$^*$ \\
\;\,$+$morph & 0.563$^*$ & 0.313$^*$ & 0.031$^*$ & 0.594$^*$ \\
\;\,$+$both & 0.625$^*$ & 0.156$^*$ & 0.031$^*$ & 0.563$^*$ \\
\midrule
\multicolumn{5}{l}{\textit{Relapse (AUROC)}} \\
$\mathcal{F}_3$      & \textbf{0.800} & 0.800 & \textbf{0.667} & 0.600 \\
\;\,$+$img           & \textbf{0.800} & 0.800 & \textbf{0.667} & 0.700 \\
\;\,$+$morph         & 0.633 & \textbf{0.833} & 0.500 & \textbf{0.800} \\
\;\,$+$both          & 0.733 & \textbf{0.833} & 0.500 & \textbf{0.800} \\
\addlinespace[0.3pt]
$\mathcal{F}_4$      & 0.600 & \textbf{0.800} & \textbf{0.567} & 0.500 \\
\;\,$+$img           & \textbf{0.800} & 0.700 & 0.500 & 0.625 \\
\;\,$+$morph         & \textbf{0.800} & 0.733 & 0.500 & \textbf{0.708} \\
\;\,$+$both          & 0.633 & 0.733 & 0.500 & 0.583 \\
\addlinespace[0.3pt]
$\mathcal{F}_5$      & \textbf{0.600} & \textbf{0.800} & 0.567 & \textbf{0.600} \\
\;\,$+$img           & \textbf{0.600} & 0.600 & \textbf{0.733} & \textbf{0.600} \\
\;\,$+$morph         & 0.433 & 0.733 & 0.500 & 0.500 \\
\;\,$+$both          & \textbf{0.600} & 0.533 & 0.400 & \textbf{0.600} \\
\addlinespace[0.3pt]
$\mathcal{F}_6$      & 0.700 & \textbf{0.900} & 0.567 & 0.500 \\
\;\,$+$img           & \textbf{0.800} & 0.800 & \textbf{0.733} & 0.600 \\
\;\,$+$morph         & 0.700 & \textbf{0.900} & 0.567 & \textbf{0.700} \\
\;\,$+$both          & \textbf{0.800} & 0.800 & 0.567 & 0.500 \\
\addlinespace[0.3pt]
$\mathcal{F}_7$      & \textbf{0.500} & \textbf{0.500} & 0.500 & \textbf{0.500} \\
\;\,$+$img           & \textbf{0.500} & \textbf{0.500} & 0.500 & \textbf{0.500} \\
\;\,$+$morph         & \textbf{0.500} & \textbf{0.500} & \textbf{0.600} & \textbf{0.500} \\
\;\,$+$both          & \textbf{0.500} & \textbf{0.500} & 0.500 & \textbf{0.500} \\
\addlinespace[2pt]
Avg & 0.640 & \textbf{0.760} & 0.573 & 0.540 \\
\;\,$+$img & \textbf{0.700} & 0.680 & \textbf{0.627} & 0.605 \\
\;\,$+$morph & 0.613 & 0.740 & 0.533 & \textbf{0.642} \\
\;\,$+$both & 0.653 & 0.680 & 0.493 & 0.597 \\
\bottomrule
\end{tabular}
\end{table}

\begin{table}[p]
\caption{\textbf{Information gain from bone marrow imaging at every
clinical feature level, external survival tasks (merged external
cohort).} Columns as in Table~\ref{tab:imgmorph_ext_cls}. C-index for
EFS and OS. Base rows are re-evaluated on the morphology-available
patient subset, so they can differ slightly from
Table~\ref{tab:phase1_ext}. Pooling the two source hospitals gives the
morphology-available subset EFS/OS events, so survival is evaluable
across the whole panel (no \textemdash). \textbf{Bold}: best variant
within each level block (per column; ties all bolded). The pooled
external fair-comparison subset is small ($n{=}13$); indicative only.}
\label{tab:imgmorph_ext_surv}
\centering
\scriptsize
\renewcommand{\arraystretch}{0.6}
\setlength{\tabcolsep}{28pt}
\begin{tabular}{@{}l cccc@{}}
\toprule
 & Gm & Gma & Iv & Qw \\
\midrule
\multicolumn{5}{l}{\textit{EFS (C-index)}} \\
$\mathcal{F}_3$      & 0.885 & 0.846 & 0.692 & \textbf{0.923} \\
\;\,$+$img           & 0.846 & 0.744 & 0.654 & 0.795 \\
\;\,$+$morph         & \textbf{0.897} & \textbf{0.910} & \textbf{0.833} & 0.744 \\
\;\,$+$both          & 0.795 & 0.872 & 0.603 & 0.782 \\
\addlinespace[2pt]
$\mathcal{F}_4$      & 0.833 & 0.846 & 0.692 & \textbf{0.923} \\
\;\,$+$img           & \textbf{0.846} & 0.744 & 0.603 & 0.782 \\
\;\,$+$morph         & 0.808 & \textbf{0.910} & \textbf{0.833} & 0.744 \\
\;\,$+$both          & 0.667 & 0.872 & 0.641 & 0.782 \\
\addlinespace[2pt]
$\mathcal{F}_5$      & 0.872 & \textbf{0.897} & 0.872 & \textbf{0.577} \\
\;\,$+$img           & \textbf{0.910} & 0.885 & 0.692 & 0.500 \\
\;\,$+$morph         & 0.859 & 0.872 & \textbf{0.885} & \textbf{0.577} \\
\;\,$+$both          & 0.769 & 0.872 & 0.731 & 0.500 \\
\addlinespace[2pt]
$\mathcal{F}_6$      & \textbf{0.885} & 0.782 & 0.859 & 0.680 \\
\;\,$+$img           & 0.846 & 0.821 & 0.756 & 0.731 \\
\;\,$+$morph         & 0.808 & \textbf{0.859} & \textbf{0.897} & \textbf{0.859} \\
\;\,$+$both          & 0.859 & 0.782 & 0.795 & 0.641 \\
\addlinespace[2pt]
$\mathcal{F}_7$      & \textbf{0.821} & 0.756 & \textbf{0.872} & 0.680 \\
\;\,$+$img           & \textbf{0.821} & 0.718 & 0.692 & 0.551 \\
\;\,$+$morph         & 0.782 & \textbf{0.782} & 0.859 & \textbf{0.692} \\
\;\,$+$both          & 0.769 & 0.718 & 0.718 & 0.590 \\
\addlinespace[2pt]
Avg & \textbf{0.859} & 0.826 & 0.797 & \textbf{0.756} \\
\;\,$+$img & 0.854 & 0.782 & 0.679 & 0.672 \\
\;\,$+$morph & 0.831 & \textbf{0.867} & \textbf{0.862} & 0.723 \\
\;\,$+$both & 0.772 & 0.823 & 0.697 & 0.659 \\
\midrule
\multicolumn{5}{l}{\textit{OS (C-index)}} \\
$\mathcal{F}_3$      & \textbf{0.850} & \textbf{0.883} & 0.600 & 0.650 \\
\;\,$+$img           & 0.817 & 0.750 & 0.600 & 0.750 \\
\;\,$+$morph         & 0.817 & 0.850 & 0.600 & \textbf{0.833} \\
\;\,$+$both          & 0.817 & 0.833 & \textbf{0.700} & 0.700 \\
\addlinespace[2pt]
$\mathcal{F}_4$      & 0.817 & \textbf{0.883} & \textbf{0.600} & 0.650 \\
\;\,$+$img           & 0.817 & 0.750 & 0.550 & 0.750 \\
\;\,$+$morph         & \textbf{0.900} & 0.850 & \textbf{0.600} & \textbf{0.833} \\
\;\,$+$both          & 0.833 & 0.833 & \textbf{0.600} & 0.700 \\
\addlinespace[2pt]
$\mathcal{F}_5$      & \textbf{0.933} & \textbf{0.933} & 0.950 & \textbf{0.850} \\
\;\,$+$img           & 0.900 & 0.783 & 0.950 & 0.500 \\
\;\,$+$morph         & 0.800 & 0.900 & 0.867 & 0.800 \\
\;\,$+$both          & 0.817 & 0.833 & \textbf{0.983} & 0.500 \\
\addlinespace[2pt]
$\mathcal{F}_6$      & 0.800 & 0.700 & 0.850 & 0.817 \\
\;\,$+$img           & 0.833 & 0.750 & 0.817 & 0.567 \\
\;\,$+$morph         & 0.817 & 0.700 & 0.833 & \textbf{0.850} \\
\;\,$+$both          & \textbf{0.850} & \textbf{0.800} & \textbf{0.917} & 0.567 \\
\addlinespace[2pt]
$\mathcal{F}_7$      & 0.817 & 0.700 & 0.833 & 0.667 \\
\;\,$+$img           & \textbf{0.850} & 0.617 & \textbf{0.933} & 0.650 \\
\;\,$+$morph         & \textbf{0.850} & 0.700 & 0.833 & \textbf{0.900} \\
\;\,$+$both          & 0.800 & \textbf{0.733} & 0.917 & 0.717 \\
\addlinespace[2pt]
Avg & \textbf{0.843} & \textbf{0.820} & 0.767 & 0.727 \\
\;\,$+$img & \textbf{0.843} & 0.730 & 0.770 & 0.643 \\
\;\,$+$morph & 0.837 & 0.800 & 0.747 & \textbf{0.843} \\
\;\,$+$both & 0.823 & 0.807 & \textbf{0.823} & 0.637 \\
\bottomrule
\end{tabular}
\end{table}

%%--------------------------------------------------------------%%
\subsection{Phase~2: Temporal validation and classification+survival
results}
\label{sec:phase2}
%%--------------------------------------------------------------%%

Phase~2 evaluates DT and RF models on three temporal splits
($\mathcal{T}_1=2016$: $n_\text{train}=99$, $n_\text{test}=201$;
$\mathcal{T}_2=2017$: $n_\text{train}=186$, $n_\text{test}=114$;
$\mathcal{T}_3=2018$: $n_\text{train}=283$, $n_\text{test}=17$)
and a random survival forest (RSF) on the same splits for the EFS
and OS targets (internal results in Table~\ref{tab:phase2_int},
external in Table~\ref{tab:phase2_ext}).

\textbf{Classification on the internal cohort (Split~2).}
Final Risk DT AUROC rose from 0.810 at $\mathcal{F}_1$ to
0.912 at $\mathcal{F}_6$, with the largest single gain at $\mathcal{F}_5$
(+0.127, introduction of Initial Risk and Day~5 blast percentage).
Including MRD at $\mathcal{F}_7$ further raised DT AUROC to
\textbf{0.949} (RF:~0.998), the highest classification result across
all conditions.
Systematically adding the SHAP-selected morphology features across
$\mathcal{F}_3$--$\mathcal{F}_7$ (Table~\ref{tab:phase2_morph}) gives
no consistent internal AUROC gain for D19MRD or the other
classification tasks; D46MRD reached DT~0.671 and RF~0.636 at
$\mathcal{F}_4$.
Relapse prediction remained difficult for Phase~2 as for Phase~1
(best DT AUROC~0.529, RF~0.630), mirroring the broader literature
on relapse biology~\cite{[CITE:Martinez2025]}.

\textbf{Survival on the internal cohort (Split~2).}
RSF at $\mathcal{F}_7$ achieved C-index~\textbf{0.746} and
time-AUC at 24~months~0.775 for EFS (log-rank $p=0.010$),
outperforming the best LLM zero-shot EFS C-index (0.716) obtained
without any training data.
For OS, RSF C-index was 0.629 ($p=0.489$; not significant), reflecting
the low event rate (8.0\%) and limited statistical power.

\begin{table}[t]
\caption{\textbf{Method comparison on the internal cohort, per temporal
split.} Rows are methods: LLM zero-shot (Gemini, best feature level,
evaluated on each split's test set); decision tree (DT, a survival tree
for EFS/OS) and random forest (RF, a random survival forest RSF for
EFS/OS), our Phase~2 models. Classification: AUROC; survival (EFS, OS):
C-index. \textbf{Bold}
marks the best method in each split. Test sets by enrolment year:
Split~1 (2017--19, $n$=201), Split~2 (2018--19, $n$=114), Split~3 (2019,
$n$=17). \textemdash{}: no OS events in the 2019 test set. Full
per-feature-level detail in Table~\ref{tab:phase2_int_full}.}
\label{tab:phase2_int}
\centering
\small
\setlength{\tabcolsep}{6pt}
\begin{tabular}{@{}l ccc@{}}
\toprule
\textbf{Method} & \textbf{Split~1} & \textbf{Split~2} & \textbf{Split~3} \\
 & \footnotesize $n$=201 & \footnotesize $n$=114 & \footnotesize $n$=17 \\
\midrule
\multicolumn{4}{l}{\textit{Final Risk (AUROC)}} \\
LLM (Gemini) & 0.860 & 0.856 & \textbf{1.000} \\
DT & 0.893 & 0.949 & 0.979 \\
RF & \textbf{0.996} & \textbf{0.998} & \textbf{1.000} \\
\multicolumn{4}{l}{\textit{D19 MRD (AUROC)}} \\
LLM (Gemini) & 0.651 & 0.745 & 0.560 \\
DT & 0.555 & 0.725 & \textbf{0.893} \\
RF & \textbf{0.668} & \textbf{0.853} & 0.833 \\
\multicolumn{4}{l}{\textit{D46 MRD (AUROC)}} \\
LLM (Gemini) & \textbf{0.614} & 0.632 & 0.560 \\
DT & 0.545 & \textbf{0.671} & 0.464 \\
RF & 0.597 & 0.636 & \textbf{0.655} \\
\multicolumn{4}{l}{\textit{Relapse (AUROC)}} \\
LLM (Gemini) & \textbf{0.593} & 0.589 & \textbf{0.786} \\
DT & 0.550 & 0.529 & 0.500 \\
RF & 0.591 & \textbf{0.630} & 0.500 \\
\multicolumn{4}{l}{\textit{EFS (C-index)}} \\
LLM (Gemini) & 0.661 & 0.674 & \textbf{0.772} \\
DT & 0.674 & 0.718 & 0.691 \\
RSF & \textbf{0.684} & \textbf{0.746} & 0.706 \\
\multicolumn{4}{l}{\textit{OS (C-index)}} \\
LLM (Gemini) & 0.682 & \textbf{0.665} & \textemdash \\
DT & 0.562 & 0.562 & \textemdash \\
RSF & \textbf{0.688} & 0.628 & \textemdash \\
\bottomrule
\end{tabular}
\end{table}

\begin{table}[p]
\caption{\textbf{Phase~2 internal results across all feature levels and
the three temporal splits} (full detail behind the summary in
Table~\ref{tab:phase2_int}). Classification: DT / RF AUROC; survival
(EFS, OS): survival-tree (DT) / RSF C-index. \textbf{Bold} marks
the best feature level in each column. Splits by enrolment year: S1
train~$\leq$16/test~17--19; S2 $\leq$17/18--19; S3 $\leq$18/19.
\textemdash{}: not evaluable (D19/D46MRD are inputs at $\mathcal{F}_7$;
no OS events in the 2019 test set).}
\label{tab:phase2_int_full}
\centering
\footnotesize
\renewcommand{\arraystretch}{0.88}
\setlength{\tabcolsep}{5pt}
\begin{tabular}{@{}l cc cc cc@{}}
\toprule
 & \multicolumn{2}{c}{\textbf{Split~1}} & \multicolumn{2}{c}{\textbf{Split~2}} & \multicolumn{2}{c}{\textbf{Split~3}} \\
\cmidrule(lr){2-3}\cmidrule(lr){4-5}\cmidrule(lr){6-7}
\textbf{Lv.} & DT & RF & DT & RF & DT & RF \\
\midrule
\multicolumn{7}{l}{\textit{Final Risk (DT / RF AUROC)}} \\
$\mathcal{F}_1$ & 0.759 & 0.842 & 0.810 & 0.871 & 0.799 & 0.806 \\
$\mathcal{F}_2$ & 0.759 & 0.820 & 0.810 & 0.830 & 0.778 & 0.847 \\
$\mathcal{F}_3$ & 0.711 & 0.849 & 0.782 & 0.872 & 0.771 & 0.743 \\
$\mathcal{F}_4$ & 0.711 & 0.852 & 0.757 & 0.847 & 0.771 & 0.799 \\
$\mathcal{F}_5$ & \textbf{0.893} & 0.897 & 0.909 & 0.943 & 0.736 & 0.847 \\
$\mathcal{F}_6$ & 0.893 & 0.922 & 0.911 & 0.951 & 0.736 & 0.882 \\
$\mathcal{F}_7$ & 0.891 & \textbf{0.996} & \textbf{0.949} & \textbf{0.998} & \textbf{0.979} & \textbf{1.000} \\
\multicolumn{7}{l}{\textit{D19 MRD (DT / RF AUROC)}} \\
$\mathcal{F}_1$ & 0.538 & 0.512 & 0.483 & 0.425 & 0.393 & 0.643 \\
$\mathcal{F}_2$ & 0.497 & 0.487 & 0.495 & 0.495 & 0.464 & 0.667 \\
$\mathcal{F}_3$ & \textbf{0.555} & 0.629 & 0.687 & 0.767 & \textbf{0.893} & \textbf{0.833} \\
$\mathcal{F}_4$ & 0.555 & 0.641 & \textbf{0.725} & \textbf{0.853} & 0.893 & 0.786 \\
$\mathcal{F}_5$ & 0.555 & \textbf{0.668} & 0.687 & 0.764 & 0.750 & 0.738 \\
$\mathcal{F}_6$ & 0.555 & 0.643 & 0.687 & 0.840 & 0.845 & 0.833 \\
$\mathcal{F}_7$ & \textemdash & \textemdash & \textemdash & \textemdash & \textemdash & \textemdash \\
\multicolumn{7}{l}{\textit{D46 MRD (DT / RF AUROC)}} \\
$\mathcal{F}_1$ & 0.364 & 0.534 & 0.405 & 0.516 & \textbf{0.464} & 0.619 \\
$\mathcal{F}_2$ & 0.519 & 0.396 & 0.405 & 0.533 & 0.321 & \textbf{0.655} \\
$\mathcal{F}_3$ & \textbf{0.545} & 0.564 & 0.519 & 0.605 & 0.369 & 0.309 \\
$\mathcal{F}_4$ & 0.531 & 0.528 & \textbf{0.671} & \textbf{0.636} & 0.369 & 0.417 \\
$\mathcal{F}_5$ & 0.545 & 0.555 & 0.548 & 0.550 & 0.369 & 0.167 \\
$\mathcal{F}_6$ & 0.545 & \textbf{0.597} & 0.548 & 0.613 & 0.369 & 0.464 \\
$\mathcal{F}_7$ & \textemdash & \textemdash & \textemdash & \textemdash & \textemdash & \textemdash \\
\multicolumn{7}{l}{\textit{Relapse (DT / RF AUROC)}} \\
$\mathcal{F}_1$ & 0.480 & 0.439 & \textbf{0.529} & 0.465 & 0.274 & 0.476 \\
$\mathcal{F}_2$ & \textbf{0.550} & 0.546 & 0.483 & 0.526 & \textbf{0.500} & 0.357 \\
$\mathcal{F}_3$ & 0.528 & 0.483 & 0.411 & 0.485 & 0.286 & 0.274 \\
$\mathcal{F}_4$ & 0.540 & 0.564 & 0.411 & 0.471 & 0.286 & 0.036 \\
$\mathcal{F}_5$ & 0.527 & 0.517 & 0.411 & 0.624 & 0.286 & 0.286 \\
$\mathcal{F}_6$ & 0.528 & 0.547 & 0.411 & \textbf{0.630} & 0.488 & \textbf{0.500} \\
$\mathcal{F}_7$ & 0.529 & \textbf{0.591} & 0.411 & 0.597 & 0.488 & 0.429 \\
\multicolumn{7}{l}{\textit{EFS (DT / RSF C-index)}} \\
$\mathcal{F}_1$ & 0.468 & 0.532 & 0.443 & 0.610 & \textbf{0.691} & \textbf{0.706} \\
$\mathcal{F}_2$ & 0.482 & 0.516 & 0.501 & 0.597 & 0.529 & 0.559 \\
$\mathcal{F}_3$ & 0.527 & 0.542 & 0.477 & 0.552 & 0.390 & 0.500 \\
$\mathcal{F}_4$ & 0.546 & 0.520 & 0.482 & 0.533 & 0.426 & 0.485 \\
$\mathcal{F}_5$ & 0.494 & 0.511 & 0.490 & 0.572 & 0.390 & 0.588 \\
$\mathcal{F}_6$ & 0.497 & 0.560 & 0.485 & 0.592 & 0.390 & 0.662 \\
$\mathcal{F}_7$ & \textbf{0.674} & \textbf{0.684} & \textbf{0.718} & \textbf{0.746} & 0.640 & 0.691 \\
\multicolumn{7}{l}{\textit{OS (DT / RSF C-index)}} \\
$\mathcal{F}_1$ & 0.541 & \textbf{0.688} & 0.531 & 0.497 & \textemdash & \textemdash \\
$\mathcal{F}_2$ & 0.467 & 0.651 & 0.541 & 0.540 & \textemdash & \textemdash \\
$\mathcal{F}_3$ & 0.512 & 0.556 & 0.440 & 0.515 & \textemdash & \textemdash \\
$\mathcal{F}_4$ & 0.498 & 0.590 & 0.472 & 0.529 & \textemdash & \textemdash \\
$\mathcal{F}_5$ & 0.492 & 0.586 & 0.445 & 0.622 & \textemdash & \textemdash \\
$\mathcal{F}_6$ & \textbf{0.562} & 0.621 & \textbf{0.562} & 0.620 & \textemdash & \textemdash \\
$\mathcal{F}_7$ & 0.562 & 0.607 & 0.505 & \textbf{0.628} & \textemdash & \textemdash \\
\bottomrule
\end{tabular}
\end{table}

\textbf{External validation (merged cohort, OOD).}
With all 301 internal patients as training, RF and DT retained strong
Final Risk discrimination on the pooled external cohort (best DT
AUROC~0.968, RF~0.991 at $\mathcal{F}_7$; Table~\ref{tab:phase2_ext}),
confirming transfer to an independent patient population.
Pooling the two source hospitals also resolves the single-class
collapse seen when they were evaluated separately: Relapse and D19MRD
are now two-class across the full external cohort and AUROC-evaluable
— best Relapse~0.583 (DT) / 0.528 (RF) and best D19MRD~0.628 (DT) /
0.635 (RF; Table~\ref{tab:phase2_ext_summary}) — while D46MRD still
has only one positive case after pooling and is reported as accuracy
(DT~0.732, RF~0.976).
For survival, the pooled cohort (24 EFS events, 7 deaths) gave best
RSF C-indices of 0.608 (EFS, $\mathcal{F}_7$) and 0.718 (OS,
$\mathcal{F}_1$/$\mathcal{F}_2$), with DT reaching 0.603 (EFS,
$\mathcal{F}_1$) and 0.710 (OS, $\mathcal{F}_7$).

\begin{table}[p]
\caption{\textbf{Phase~2 out-of-distribution results on the pooled
external cohort, across all feature levels} (train~=~all 301 internal
patients; test~=~the two source hospitals pooled into a single
external cohort, $n$=41). Classification: DT / RF AUROC; $^*$~=
accuracy (D46MRD has only one positive case after pooling, so accuracy
is reported instead). Survival (EFS, OS): survival-tree (DT) /
random-survival-forest (RSF) C-index. \textbf{Bold} marks the best
feature level in each column. \textemdash{}: not evaluable (D19/D46MRD
are inputs at $\mathcal{F}_7$). Pooled external cohort; small, indicative
only.}
\label{tab:phase2_ext}
\centering
\footnotesize
\renewcommand{\arraystretch}{0.84}
\setlength{\tabcolsep}{5pt}
\begin{tabular}{@{}l cc@{}}
\toprule
\textbf{Lv.} & DT & RF \\
\midrule
\multicolumn{3}{l}{\textit{Final Risk (DT / RF)}} \\
$\mathcal{F}_1$ & 0.805 & 0.878 \\
$\mathcal{F}_2$ & 0.827 & 0.863 \\
$\mathcal{F}_3$ & 0.779 & 0.843 \\
$\mathcal{F}_4$ & 0.778 & 0.750 \\
$\mathcal{F}_5$ & 0.912 & 0.938 \\
$\mathcal{F}_6$ & 0.912 & 0.934 \\
$\mathcal{F}_7$ & \textbf{0.968} & \textbf{0.991} \\
\multicolumn{3}{l}{\textit{D19 MRD (DT / RF)}} \\
$\mathcal{F}_1$ & \textbf{0.628} & 0.591 \\
$\mathcal{F}_2$ & 0.466 & \textbf{0.635} \\
$\mathcal{F}_3$ & 0.382 & 0.335 \\
$\mathcal{F}_4$ & 0.382 & 0.341 \\
$\mathcal{F}_5$ & 0.392 & 0.307 \\
$\mathcal{F}_6$ & 0.392 & 0.453 \\
$\mathcal{F}_7$ & \textemdash & \textemdash \\
\multicolumn{3}{l}{\textit{D46 MRD (DT / RF)}} \\
$\mathcal{F}_1$ & \textbf{0.732$^*$} & \textbf{0.976$^*$} \\
$\mathcal{F}_2$ & 0.561$^*$ & 0.976$^*$ \\
$\mathcal{F}_3$ & 0.488$^*$ & 0.976$^*$ \\
$\mathcal{F}_4$ & 0.244$^*$ & 0.976$^*$ \\
$\mathcal{F}_5$ & 0.488$^*$ & 0.976$^*$ \\
$\mathcal{F}_6$ & 0.488$^*$ & 0.976$^*$ \\
$\mathcal{F}_7$ & \textemdash & \textemdash \\
\multicolumn{3}{l}{\textit{Relapse (DT / RF)}} \\
$\mathcal{F}_1$ & \textbf{0.583} & \textbf{0.528} \\
$\mathcal{F}_2$ & 0.397 & 0.494 \\
$\mathcal{F}_3$ & 0.333 & 0.300 \\
$\mathcal{F}_4$ & 0.525 & 0.272 \\
$\mathcal{F}_5$ & 0.333 & 0.377 \\
$\mathcal{F}_6$ & 0.219 & 0.267 \\
$\mathcal{F}_7$ & 0.250 & 0.383 \\
\multicolumn{3}{l}{\textit{EFS (DT / RSF C-index)}} \\
$\mathcal{F}_1$ & \textbf{0.603} & 0.567 \\
$\mathcal{F}_2$ & 0.552 & 0.572 \\
$\mathcal{F}_3$ & 0.507 & 0.434 \\
$\mathcal{F}_4$ & 0.509 & 0.475 \\
$\mathcal{F}_5$ & 0.419 & 0.444 \\
$\mathcal{F}_6$ & 0.419 & 0.439 \\
$\mathcal{F}_7$ & 0.490 & \textbf{0.608} \\
\multicolumn{3}{l}{\textit{OS (DT / RSF C-index)}} \\
$\mathcal{F}_1$ & 0.619 & \textbf{0.718} \\
$\mathcal{F}_2$ & 0.657 & 0.718 \\
$\mathcal{F}_3$ & 0.571 & 0.468 \\
$\mathcal{F}_4$ & 0.702 & 0.512 \\
$\mathcal{F}_5$ & 0.647 & 0.556 \\
$\mathcal{F}_6$ & 0.647 & 0.512 \\
$\mathcal{F}_7$ & \textbf{0.710} & 0.643 \\
\bottomrule
\end{tabular}
\end{table}

\begin{table}[t]
\caption{\textbf{Method comparison on the pooled external cohort}
(out-of-distribution; no temporal split; the two source hospitals
pooled into a single cohort, $n$=41). Rows are methods: LLM zero-shot
(Gemini, best feature level); decision tree (DT, a survival tree for
EFS/OS) and random forest (RF, RSF for EFS/OS), trained on all 301
internal patients. Classification: AUROC ($^*$~= accuracy; D46MRD has
only one positive case after pooling, so accuracy is reported instead);
survival (EFS, OS): C-index. \textbf{Bold} marks the best method for
each task. Pooled external cohort; small, indicative only. Full
per-feature-level detail in Table~\ref{tab:phase2_ext}.}
\label{tab:phase2_ext_summary}
\centering
\small
\setlength{\tabcolsep}{7pt}
\begin{tabular}{@{}l c@{}}
\toprule
\textbf{Method} & \textbf{External} \\
 & \footnotesize $n$=41 \\
\midrule
\multicolumn{2}{l}{\textit{Final Risk (AUROC)}} \\
LLM (Gemini) & 0.925 \\
DT & 0.968 \\
RF & \textbf{0.991} \\
\multicolumn{2}{l}{\textit{D19 MRD (AUROC)}} \\
LLM (Gemini) & 0.550 \\
DT & 0.628 \\
RF & \textbf{0.635} \\
\multicolumn{2}{l}{\textit{D46 MRD (accuracy)}} \\
LLM (Gemini) & 0.690$^*$ \\
DT & 0.732$^*$ \\
RF & \textbf{0.976$^*$} \\
\multicolumn{2}{l}{\textit{Relapse (AUROC)}} \\
LLM (Gemini) & \textbf{0.717} \\
DT & 0.583 \\
RF & 0.528 \\
\multicolumn{2}{l}{\textit{EFS (C-index)}} \\
LLM (Gemini) & \textbf{0.715} \\
DT & 0.603 \\
RSF & 0.608 \\
\multicolumn{2}{l}{\textit{OS (C-index)}} \\
LLM (Gemini) & \textbf{0.810} \\
DT & 0.710 \\
RSF & 0.718 \\
\bottomrule
\end{tabular}
\end{table}

\begin{table}[p]
\caption{\textbf{Morphology information gain for the Phase~2 random forest
across clinical feature levels.} For each base level $\mathcal{F}_k$
($k=3\text{--}7$), the $+$morph row adds the SHAP-selected quantitative
bone-marrow morphology features (Section~\ref{sec:imgmorph}); raw images
cannot enter a tabular RF. Classification: RF AUROC; EFS/OS: RSF C-index.
Columns are the three internal temporal splits and the pooled external
cohort ($n$=41, indicative only). \textbf{Bold}: $+$morph exceeds the
same-level base in that column. \textemdash: OS is not evaluable in the
2019 test split (Split~3), which has no observed OS events. D19/D46MRD
are shown only up to $\mathcal{F}_6$, as MRD status is an input feature
at $\mathcal{F}_7$.}
\label{tab:phase2_morph}
\centering
\footnotesize
\setlength{\tabcolsep}{5pt}
\begin{tabular}{@{}ll cccc@{}}
\toprule
\textbf{Lv.} & & \textbf{Split~1} & \textbf{Split~2} & \textbf{Split~3} & \textbf{Ext} \\
\midrule
\multicolumn{6}{l}{\textit{Final Risk (RF AUROC)}} \\
$\mathcal{F}_3$ & base & 0.849 & 0.872 & 0.743 & 0.843 \\
& \;\,$+$morph & \textbf{0.859} & 0.868 & \textbf{0.854} & 0.799 \\
\addlinespace[2pt]
$\mathcal{F}_4$ & base & 0.852 & 0.847 & 0.799 & 0.750 \\
& \;\,$+$morph & \textbf{0.874} & \textbf{0.859} & 0.729 & \textbf{0.808} \\
\addlinespace[2pt]
$\mathcal{F}_5$ & base & 0.897 & 0.943 & 0.847 & 0.938 \\
& \;\,$+$morph & \textbf{0.916} & 0.932 & \textbf{0.882} & 0.913 \\
\addlinespace[2pt]
$\mathcal{F}_6$ & base & 0.922 & 0.951 & 0.882 & 0.934 \\
& \;\,$+$morph & \textbf{0.926} & 0.938 & 0.882 & 0.914 \\
\addlinespace[2pt]
$\mathcal{F}_7$ & base & 0.996 & 0.998 & 1.000 & 0.991 \\
& \;\,$+$morph & 0.994 & \textbf{0.999} & 1.000 & 0.976 \\
\midrule
\multicolumn{6}{l}{\textit{D19 MRD (RF AUROC)}} \\
$\mathcal{F}_3$ & base & 0.629 & 0.767 & 0.833 & 0.335 \\
& \;\,$+$morph & \textbf{0.704} & 0.754 & \textbf{0.881} & 0.294 \\
\addlinespace[2pt]
$\mathcal{F}_4$ & base & 0.641 & 0.853 & 0.786 & 0.341 \\
& \;\,$+$morph & \textbf{0.709} & 0.841 & \textbf{0.833} & 0.277 \\
\addlinespace[2pt]
$\mathcal{F}_5$ & base & 0.668 & 0.764 & 0.738 & 0.307 \\
& \;\,$+$morph & \textbf{0.689} & \textbf{0.785} & \textbf{0.881} & 0.307 \\
\addlinespace[2pt]
$\mathcal{F}_6$ & base & 0.643 & 0.840 & 0.833 & 0.453 \\
& \;\,$+$morph & \textbf{0.712} & 0.766 & 0.762 & 0.182 \\
\midrule
\multicolumn{6}{l}{\textit{D46 MRD (RF AUROC)}} \\
$\mathcal{F}_3$ & base & 0.564 & 0.605 & 0.310 & 0.975 \\
& \;\,$+$morph & \textbf{0.606} & 0.470 & \textbf{0.393} & 0.925 \\
\addlinespace[2pt]
$\mathcal{F}_4$ & base & 0.528 & 0.636 & 0.417 & 0.725 \\
& \;\,$+$morph & 0.527 & 0.595 & 0.286 & \textbf{0.925} \\
\addlinespace[2pt]
$\mathcal{F}_5$ & base & 0.555 & 0.550 & 0.167 & 0.975 \\
& \;\,$+$morph & \textbf{0.609} & \textbf{0.558} & \textbf{0.452} & \textbf{0.988} \\
\addlinespace[2pt]
$\mathcal{F}_6$ & base & 0.597 & 0.613 & 0.464 & 1.000 \\
& \;\,$+$morph & 0.576 & 0.468 & \textbf{0.595} & 0.938 \\
\midrule
\multicolumn{6}{l}{\textit{Relapse (RF AUROC)}} \\
$\mathcal{F}_3$ & base & 0.483 & 0.485 & 0.274 & 0.300 \\
& \;\,$+$morph & \textbf{0.521} & \textbf{0.497} & 0.143 & \textbf{0.390} \\
\addlinespace[2pt]
$\mathcal{F}_4$ & base & 0.564 & 0.471 & 0.036 & 0.272 \\
& \;\,$+$morph & 0.530 & \textbf{0.506} & \textbf{0.167} & \textbf{0.337} \\
\addlinespace[2pt]
$\mathcal{F}_5$ & base & 0.517 & 0.624 & 0.286 & 0.377 \\
& \;\,$+$morph & \textbf{0.533} & 0.522 & \textbf{0.333} & 0.373 \\
\addlinespace[2pt]
$\mathcal{F}_6$ & base & 0.547 & 0.630 & 0.500 & 0.267 \\
& \;\,$+$morph & \textbf{0.556} & 0.519 & 0.345 & \textbf{0.333} \\
\addlinespace[2pt]
$\mathcal{F}_7$ & base & 0.591 & 0.597 & 0.429 & 0.383 \\
& \;\,$+$morph & 0.491 & 0.430 & 0.214 & \textbf{0.449} \\
\midrule
\multicolumn{6}{l}{\textit{EFS (RSF C-index)}} \\
$\mathcal{F}_3$ & base & 0.542 & 0.552 & 0.500 & 0.434 \\
& \;\,$+$morph & 0.530 & 0.516 & 0.397 & \textbf{0.439} \\
\addlinespace[2pt]
$\mathcal{F}_4$ & base & 0.520 & 0.533 & 0.485 & 0.475 \\
& \;\,$+$morph & 0.517 & \textbf{0.559} & 0.368 & 0.455 \\
\addlinespace[2pt]
$\mathcal{F}_5$ & base & 0.511 & 0.572 & 0.588 & 0.444 \\
& \;\,$+$morph & \textbf{0.547} & 0.520 & 0.471 & 0.435 \\
\addlinespace[2pt]
$\mathcal{F}_6$ & base & 0.560 & 0.592 & 0.662 & 0.439 \\
& \;\,$+$morph & 0.549 & 0.556 & 0.353 & \textbf{0.477} \\
\addlinespace[2pt]
$\mathcal{F}_7$ & base & 0.684 & 0.746 & 0.691 & 0.608 \\
& \;\,$+$morph & 0.668 & 0.671 & 0.603 & \textbf{0.626} \\
\midrule
\multicolumn{6}{l}{\textit{OS (RSF C-index)}} \\
$\mathcal{F}_3$ & base & 0.556 & 0.515 & \textemdash & 0.468 \\
& \;\,$+$morph & 0.536 & 0.415 & \textemdash & \textbf{0.659} \\
\addlinespace[2pt]
$\mathcal{F}_4$ & base & 0.590 & 0.529 & \textemdash & 0.512 \\
& \;\,$+$morph & 0.541 & 0.508 & \textemdash & \textbf{0.714} \\
\addlinespace[2pt]
$\mathcal{F}_5$ & base & 0.586 & 0.622 & \textemdash & 0.556 \\
& \;\,$+$morph & 0.570 & 0.542 & \textemdash & \textbf{0.718} \\
\addlinespace[2pt]
$\mathcal{F}_6$ & base & 0.621 & 0.620 & \textemdash & 0.512 \\
& \;\,$+$morph & 0.611 & 0.550 & \textemdash & \textbf{0.762} \\
\addlinespace[2pt]
$\mathcal{F}_7$ & base & 0.607 & 0.629 & \textemdash & 0.643 \\
& \;\,$+$morph & 0.599 & 0.618 & \textemdash & \textbf{0.671} \\
\bottomrule
\end{tabular}
\end{table}

%%--------------------------------------------------------------%%
\subsection{Cosine RAG belief revision}
\label{sec:rag}
%%--------------------------------------------------------------%%

For each test patient $i$, the $K=5$ most similar training cases were
retrieved by cosine similarity in the standardised feature space and
used to compute a similarity-weighted RAG evidence probability
$p_\text{rag}$ (Eq.~\ref{eq:prag}).
The revised probability is
$p_\text{rev} = \alpha\,p_0 + (1-\alpha)\,p_\text{rag}$, where
$\alpha = 1-(1-\alpha_\text{base})\sqrt{\bar{s}}$ and
$\alpha_\text{base}=0.6$.

\textbf{Internal cohort.}
Table~\ref{tab:belief_revision} reports AUROC before and after belief
revision.
Belief revision improved performance on Final Risk (Split~2:
$+0.002$), D19MRD (Split~2: $+0.012$; Split~3: $+0.024$), and most
substantially on D46MRD (Split~2: $+0.046$; Split~3: $0.655 \to 0.881$,
$\Delta=+0.226$).
The large D46MRD gain on Split~3 reflects high retrieval quality: with
283 training cases, the five nearest neighbours share strong
phenotypic similarity (mean cosine $\bar{s} \approx 0.91$) and
concordant D46MRD outcomes.
Belief revision degraded Relapse performance (Split~2:
$0.630 \to 0.588$, $\Delta=-0.042$): because relapse prediction is
intrinsically noisy, retrieved neighbours have mixed outcomes, and
RAG evidence acts as a noise source rather than a corrective signal.

\textbf{External (pooled across sites).}
On the pooled external cohort, belief revision improved Final Risk
($0.991 \to 1.000$, $+0.009$).
D19MRD AUROCs degraded ($-0.064$) because the external set has very
few D19MRD-positive patients ($n_\text{pos}=4$), making retrieved
neighbour outcomes uninformative.
D46MRD belief revision could not be computed ($n_\text{pos}=1$).
These results confirm the internal boundary condition: belief revision
benefits tasks where the feature space clusters meaningfully by label,
but is counter-productive when class representation is too sparse to
yield informative neighbours.

\textbf{Survival belief revision.}
Extending belief revision to EFS and OS requires handling censored
neighbour outcomes.
A naive adaptation that treats the event indicator as a binary outcome
(ignoring censoring) is near-neutral to slightly negative
(EFS mean $\Delta$C-index $-0.003$, OS $-0.007$).
Replacing that evidence with a censoring-aware similarity-weighted
Kaplan--Meier neighbourhood event probability at 24 months
(Table~\ref{tab:belief_revision}, survival rows) instead applies the
revision to the normalised RSF risk score and yields small but
consistent gains: mean $\Delta$C-index $+0.006$ over all 70
experiment$\times$split$\times$endpoint cells (Wilcoxon $p=0.001$),
significant for internal EFS ($+0.006$, $p=0.004$), internal OS
($+0.006$, $p=0.016$) and external OS ($+0.014$, $p=0.045$); external
EFS was neutral.
The gains are modest and concentrated in mid-ranked configurations—the
best RSF models are already near-optimal and change little—consistent
with belief revision helping most where the base model is imperfect.

\begin{table}[ht]
\caption{Cosine RAG belief revision: RF AUROC before ($p_\text{RF}$) and
after ($p_\text{rev}$) belief revision ($K=5$, $\alpha_\text{base}=0.6$),
best experiment per task. Internal results are reported per temporal
split; external results use the pooled external cohort, as the per-site
external cohorts are too small or single-class for this analysis.
For the survival endpoints (EFS, OS) the values are Harrell's C-index:
base is the RSF risk score and the revision replaces the
similarity-weighted binary-outcome mean by a similarity-weighted
Kaplan--Meier neighbourhood event probability at 24 months
(experiment~P2-Exp8, the MRD-augmented feature set).
\textbf{Bold}: largest gain. $\dagger$~insufficient positives; not
computed. \textemdash: too few events (OS, 2019 test split).}
\label{tab:belief_revision}
\centering
\small
\setlength{\tabcolsep}{7pt}
\begin{tabular}{@{}ll ccc@{}}
\toprule
\textbf{Task} & \textbf{Split} & $p_\text{RF}$ & $p_\text{rev}$ & \textbf{$\Delta$} \\
\midrule
\multicolumn{5}{l}{\textit{Internal (temporal splits)}} \\
Final Risk & Split~1 & 0.996 & 0.998 & $+0.002$ \\
Final Risk & Split~2 & 0.998 & 0.999 & $+0.002$ \\
Final Risk & Split~3 & 1.000 & 1.000 & $+0.000$ \\
D19MRD & Split~1 & 0.629 & 0.631 & $+0.002$ \\
D19MRD & Split~2 & 0.767 & 0.779 & $+0.012$ \\
D19MRD & Split~3 & 0.833 & 0.857 & $+0.024$ \\
D46MRD & Split~1 & 0.396 & 0.371 & $-0.025$ \\
D46MRD & Split~2 & 0.533 & 0.579 & $+0.046$ \\
D46MRD & Split~3 & 0.655 & \textbf{0.881} & $\mathbf{+0.226}$ \\
Relapse & Split~1 & 0.547 & 0.546 & $-0.001$ \\
Relapse & Split~2 & 0.630 & 0.588 & $-0.042$ \\
Relapse & Split~3 & 0.500 & 0.476 & $-0.024$ \\
\addlinespace[2pt]
\multicolumn{5}{l}{\quad\textit{Survival endpoints (C-index; RSF base vs KM revision)}} \\
EFS & Split~1 & 0.684 & 0.704 & $+0.020$ \\
EFS & Split~2 & 0.746 & 0.745 & $-0.001$ \\
EFS & Split~3 & 0.691 & 0.706 & $+0.015$ \\
OS & Split~1 & 0.607 & 0.617 & $+0.010$ \\
OS & Split~2 & 0.629 & 0.634 & $+0.005$ \\
OS & Split~3 & \multicolumn{3}{c}{\textemdash} \\
\midrule
\multicolumn{5}{l}{\textit{External (pooled cohort)}} \\
Final Risk & --- & 0.991 & 1.000 & $+0.009$ \\
D19MRD & --- & 0.335 & 0.270 & $-0.064$ \\
D46MRD & --- & \multicolumn{3}{c}{$\dagger$ ($n_\text{pos}=1$)} \\
Relapse & --- & 0.267 & 0.283 & $+0.016$ \\
\addlinespace[2pt]
\multicolumn{5}{l}{\quad\textit{Survival endpoints (C-index; RSF base vs KM revision)}} \\
EFS & --- & 0.608 & 0.613 & $+0.005$ \\
OS & --- & 0.643 & 0.659 & $+0.016$ \\
\bottomrule
\end{tabular}
\end{table}

%%--------------------------------------------------------------%%
\subsection{LLM clinical narratives and qualitative evaluation}
\label{sec:narratives}
%%--------------------------------------------------------------%%

After belief revision, an LLM prompt was constructed for each test
patient containing: (i) the patient's serialised feature vector;
(ii) the RF-predicted label, revised probability $p_\text{rev}$, and
the human-readable DT decision path from root to leaf;
(iii) the six patient-salient features ranked by
$\text{importance}_j \times |z_{ij}|$;
and (iv) the structured profiles and outcomes of the top-5 retrieved
historical cases.
Gemini-3.1-Flash-Lite produced 3--5-sentence narratives for each
prediction; all narratives from the $\mathcal{F}_7$/Split~3 experiment
($n=17$ test patients) are available for physician review.

Two representative cases from the $\mathcal{F}_7$/Split~3 experiment
illustrate the range of narrative quality (Fig.~\ref{fig:rag}).
In Case~1 (correctly classified IR, persistent MRD), the top-5
retrieved neighbours all share elevated MRD and cytogenetic
abnormalities, and the narrative coherently attributes the high-risk
classification to MRD non-clearance corroborated by historical
precedent.
In Case~2 (borderline case), the narrative highlights conflicting
signals between MRD clearance and cytogenetic risk, suggesting
next-generation sequencing as a way to resolve ambiguity.

Quantitative physician evaluation of narrative quality is planned
and will follow the protocol in Table~\ref{tab:narrative_eval}.
Evaluation will cover: factual consistency with the patient's features,
relevance of historical case comparisons, clinical plausibility of the
interpretation, and overall helpfulness for clinical decision-making.

\begin{table}[ht]
\caption{Planned physician evaluation of LLM clinical narratives.
Each narrative will be assessed on four dimensions by two independent
clinicians using a 5-point Likert scale.
Results to be filled following structured review.}
\label{tab:narrative_eval}
\centering
\small
\begin{tabular}{lccccc}
\toprule
\textbf{Case ID}
  & \textbf{Factual accuracy}
  & \textbf{Retrieval relevance}
  & \textbf{Clinical plausibility}
  & \textbf{Decision helpfulness}
  & \textbf{Overall} \\
  & (1--5) & (1--5) & (1--5) & (1--5) & (1--5) \\
\midrule
Case 001 & — & — & — & — & — \\
Case 002 & — & — & — & — & — \\
Case 003 & — & — & — & — & — \\
$\vdots$ & $\vdots$ & $\vdots$ & $\vdots$ & $\vdots$ & $\vdots$ \\
Case 017 & — & — & — & — & — \\
\midrule
Mean (SD) & — & — & — & — & — \\
Inter-rater $\kappa$ & — & — & — & — & — \\
\bottomrule
\end{tabular}
\end{table}

%%==============================================================%%
\section{Discussion}\label{sec:discussion}
%%==============================================================%%

This study presents a temporally-structured multimodal framework
evaluated jointly across six ALL prediction targets—four classification
tasks and two survival endpoints—and validated at an independent
external centre.
Several findings warrant discussion.

\textbf{MRD as the dominant signal across all targets.}
The largest single performance gain in Phase~2—DT AUROC for Final Risk
from 0.912 ($\mathcal{F}_6$, cytogenetics) to 0.949 ($\mathcal{F}_7$,
plus D19/D46 MRD)—quantitatively confirms the central role of MRD
in clinical risk stratification~\cite{[CITE:Conter2010]}.
The same pattern holds for LLM EFS C-index, which peaked at feature
levels including MRD (Gemma: 0.716 at $\mathcal{F}_7$), and for
Phase~2 RSF EFS (C-index~0.746, $p=0.010$).
These convergent results from four LLMs, decision trees, and a survival
forest provide multi-method evidence that residual disease burden at
Day~46 is the strongest prognostic signal accessible before relapse.

\textbf{LLM zero-shot as a competitive baseline.}
Phase~1 Final Risk AUROCs of 0.848--0.891 without labelled training
data are competitive with Phase~2 DT at $\mathcal{F}_6$ (0.912) and
only modestly below the Phase~2 DT peak (0.949 at $\mathcal{F}_7$).
This indicates that clinical knowledge encoded during LLM pretraining
includes a meaningful understanding of paediatric ALL risk factors, and
that LLMs may serve as strong priors when historical case databases are
small or absent—for example, at the time a new treatment protocol is
introduced.
On the merged external cohort, LLM Final Risk AUROCs at the best
feature level ($\mathcal{F}_7$) ranged 0.750--0.925 across models,
surpassing the internal zero-shot values for some models, likely due
to differences in class distribution (larger IR/HR proportion in the
external cohort) rather than genuine improvement in discriminative
ability.

\textbf{Imaging information: morphology text, not raw pixels.}
The multimodal comparison (Section~\ref{sec:imgmorph}) shows an
asymmetry at diagnosis: the structured morphology profile improved
zero-shot relapse prediction for all four LLMs ($\Delta$AUROC
$+0.003$ to $+0.040$), whereas raw smear images improved Final Risk
for three of four models only marginally and degraded relapse
prediction for three of four.
This suggests that current multimodal LLMs extract prognostic
morphology more reliably from a quantitative, pipeline-derived text
profile than from raw microscopy, supporting automated morphometry as
the practical route for injecting imaging information into clinical
risk models.
The completed comparisons across all five feature levels reinforce
this asymmetry: the morphology profile improved relapse prediction for
a majority of models at three of five levels and EFS in 15 of 20
model--level comparisons, and even at $\mathcal{F}_7$—with MRD already
among the inputs—it improved Final Risk for all four models ($+0.005$
to $+0.102$), whereas raw images degraded Final Risk there for all
four.
The effects remain small and unadjusted, but they indicate that the
morphometric text profile retains additive prognostic value even after
molecular and MRD data arrive, while raw pixels do not.

\textbf{External validation: generalisation and limitations.}
Phase~2 RF Final Risk AUROC of 0.991 on the merged external cohort
(versus 0.998 internally) demonstrates strong generalisation across
independent institutions.
However, D19MRD AUROC dropped from 0.853 (internal RF at
$\mathcal{F}_4$, Split~2, Table~\ref{tab:phase2_int_full}) to 0.635
(external), indicating feature-distribution shift that partly limits
generalisation for this target.
The single D46MRD-positive patient remaining after pooling ($n=1$ of
29 labelled patients) precluded belief revision and produced unstable
AUC estimates, highlighting the need for larger external cohorts with
balanced outcome distributions.
The pooled external cohort remains small (41 patients from two
hospitals of 20 and 21) and includes a relapse-enriched referral
sub-cohort (see Limitations), so some targets remain
event-sparse even after pooling (e.g., D46MRD, $n=1$ positive);
larger, prospectively balanced external cohorts remain necessary for
a definitive test of cross-institution robustness.

\textbf{Integrating six targets in a unified framework.}
Treating EFS and OS as targets~5--6 of the same feature hierarchy—rather
than studying them in a separate survival analysis—reveals a coherent
picture: the same feature level ($\mathcal{F}_7$) that maximises
Final Risk classification also maximises EFS survival discrimination
(RSF C-index~0.746), while Relapse and OS remain challenging regardless
of feature set.
This joint view avoids the selective reporting that can arise when
survival and classification tasks are presented in separate papers or
using different patient subsets.

\textbf{Boundary conditions for RAG belief revision.}
Belief revision delivers quantitative gains specifically when the feature
space clusters meaningfully by label.
For D46MRD at Split~3 (mean similarity $\bar{s}\approx 0.91$), the
$\Delta+0.226$ AUROC gain is substantial and clinically relevant.
For Relapse ($\Delta-0.042$ internally, $+0.016$ externally with very
low base AUROC), and for OS (mean $\Delta$C-index $=-0.007$), retrieved
neighbours are too heterogeneous to provide reliable evidence.
A practical deployment criterion would be to apply belief revision
only when inter-neighbour outcome agreement exceeds a calibrated
threshold—for example, Fleiss $\kappa > 0.4$ over the top-$K$
retrieved labels.

\textbf{Relapse remains an open problem.}
Relapse prediction was intractable for all Phase~1 and Phase~2 models
(best DT AUROC~0.529, LLM best~0.643), consistent with the broader
literature~\cite{[CITE:Martinez2025],[CITE:Good2018]}.
Our feature set does not include genomic or epigenomic data (e.g.,
copy-number variation, DNA methylation), which have shown modest
incremental utility~\cite{[CITE:Mosquera2024]}.
Future multi-omics integration within the temporal feature hierarchy
may improve relapse prediction.

\textbf{Limitations.}
The internal cohort is a single-institution retrospective study, and
the external cohort—though pooled across two hospitals—remains small;
multi-centre prospective validation at greater scale is required
before clinical deployment.
The internal cohort ($n=301$) limits statistical power for OS
(24 events) and Split~3 ($n_\text{test}=17$).
The imaging comparison is complete at all feature levels internally,
but its external fair-comparison subsets remain small (classification
$n=8$, survival $n=13$); pooling now makes external survival
evaluable across the full feature range, but the imaging findings
should still be read primarily as internal-cohort evidence.
The external cohort pools two hospitals, one of which (Fudan) is a
relapse-enriched referral cohort in which all patients relapsed;
the pooled external relapse rate is therefore far above that of an
unselected ALL population, so external metrics—particularly for
relapse—are exploratory rather than representative estimates;
moreover, the complete-label subset used for the Phase~1
classification metrics ($n{=}29$) is more skewed still (23 of 29
relapsed, 79\%), because joint-label completeness retained all
relapse-positive patients.
Phase~1 LLMs were evaluated zero-shot; fine-tuning on labelled clinical
data may yield further improvements.
Physician evaluation of LLM narratives (Table~\ref{tab:narrative_eval})
is ongoing and will determine whether the narratives are clinically
actionable beyond proof-of-concept case studies.

%%==============================================================%%
\section{Methods}\label{sec:methods}
%%==============================================================%%

\subsection{Datasets and ethics}\label{sec:dataset}

The internal training cohort comprised $n=301$ patients diagnosed with
ALL at [HOSPITAL NAME] between 2015 and 2019.
All patients were treated with risk-stratified chemotherapy protocols
consistent with contemporary Chinese paediatric oncology guidelines.
Ethics approval was obtained from [ETHICS NUMBER]; all data were
de-identified prior to analysis.
Each patient record includes structured clinical data (demographics,
haematological indices, immunophenotype, 51-marker flow cytometry panel,
cytogenetics, MRD values) and bone marrow aspirate microscopy images
(Giemsa-stained, $\sim$20 fields per patient).

The external validation cohort ($n=41$) was supplied as a single
merged file spanning two source hospitals, distinguished by the
recorded hospital field: Lianyungang First People's Hospital (20
patients, 2019--2023) and Children's Hospital of Fudan University,
Shanghai (21 patients, 2017--2021, a relapse-enriched referral cohort
in which all patients relapsed); ethics approval: [ETHICS NUMBERS,
external cohort]. We pool the two hospitals and evaluate the external
cohort as a single group throughout (Section~\ref{sec:ood}).
No external data were used in any model training.

\subsection{Temporally-structured feature hierarchy}\label{sec:features}

Feature sets are temporally ordered and, except for the
$\mathcal{F}_4$ branch, form a strictly nested sequence
$\mathcal{F}_1 \subseteq \mathcal{F}_2 \subseteq \mathcal{F}_3
\subseteq \mathcal{F}_5 \subseteq \mathcal{F}_6 \subseteq \mathcal{F}_{7}$
($\mathcal{F}_4$ branches from $\mathcal{F}_3$; Table~\ref{tab:featurelist}),
where each level contains exactly the features available at the
corresponding clinical timepoint.
Boundaries: $\mathcal{F}_1$~(Day~0): WBC count and age;
$\mathcal{F}_2$~(Day~1): adds blast percentage;
$\mathcal{F}_3$~(Day~1): adds immunophenotype and 51 flow cytometry markers;
$\mathcal{F}_4$~(Day~1): adds non-risk clinical factors;
$\mathcal{F}_5$~(Day~5): adds CNS status, Day~5 blast, and Initial Risk;
$\mathcal{F}_6$~(Day~7--10): adds Fusion gene, Karyotype, FISH;
$\mathcal{F}_7$~(Day~46): adds D19MRD and D46MRD.
The complete variable list for every level is given in
Table~\ref{tab:featurelist}.
All six prediction targets are evaluated on the same feature hierarchy.

For Phase~2, categorical features were one-hot encoded and numerical
features z-score normalised using training-split statistics only:
\begin{equation}
z_{ij}^{(\text{test})} =
\frac{x_{ij}^{(\text{test})} - \mu_j^{(\text{train})}}
     {\sigma_j^{(\text{train})}},
\label{eq:zscore}
\end{equation}
with missing values imputed by training-split medians.

\subsection{Bone marrow morphology pipeline}\label{sec:morph}

\textbf{Stage~1: Cell segmentation.}
SAM~3.1~\cite{[CITE:SAM]} was applied with the text prompt
\textit{``complete dark-stained nucleated cell''} after colour-based
pre-filtering (HSV hue 120--160°) to restrict sampling to
haematoxylin-stained regions.

\textbf{Stage~2: Morphological feature extraction.}
For each retained cell, $\sim$37 features were computed: geometric
descriptors (area, perimeter, circularity, eccentricity, aspect ratio),
textural descriptors (GLCM contrast and correlation), nucleus-to-cytoplasm
(N/C) ratio, nucleoli count, and vacuole percentage.

\textbf{Stage~3: Cell-type clustering.}
BiomedCLIP~\cite{[CITE:BiomedCLIP]} was applied in a zero-shot
clustering mode with $k=7$ prototype classes (blast, granulocyte,
megakaryocyte, monocyte, plasma cell, red cell, other), with Hungarian
algorithm assignment.

\textbf{Stage~4: Per-patient aggregation.}
Cell-level features were aggregated into per-patient morphology
statistics (cell-type proportions and summary descriptors) that augment
the tabular Phase~2 models; the same profile is serialised as the
$+$morph text input in the multimodal comparison
(Section~\ref{sec:imgmorph}).

\subsection{Phase~1: LLM zero-shot inference}\label{sec:llm_inference}

For classification tasks, each patient's features were serialised into a
key--value list and presented with a task-specific system prompt
requesting a bracketed label (e.g., \texttt{[[IR]]}/\texttt{[[LR]]}).
For survival risk scoring, the same serialised features were presented
with a prompt requesting a continuous event probability (0.0--1.0).
All experiments used temperature~$=0$ (greedy decoding).
Per-patient outputs were cached; the cohort was scored for every
combination of feature level ($\times 7$, plus the imaging variants of
Section~\ref{sec:imgmorph}), task ($\times 6$), and model ($\times 4$).

\subsection{Phase~2: Decision tree and random forest}\label{sec:dt}

A decision tree (DT; \texttt{max\_depth}=4, \texttt{class\_weight}=%
`balanced', \texttt{criterion}=`gini', \texttt{random\_state}=42)
and a random forest (RF; $n_\text{estimators}=100$,
\texttt{class\_weight}=`balanced') were trained for each feature level
and temporal split.
Three temporal splits defined by enrolment year threshold $\mathcal{T}_s$
(Eq.~\ref{eq:splits}) ensure strict temporal non-leakage.
\begin{equation}
\mathcal{D}_s^{\text{train}} = \{i : \tau_i \leq \mathcal{T}_s\},
\quad
\mathcal{D}_s^{\text{test}} = \{i : \tau_i > \mathcal{T}_s\}.
\label{eq:splits}
\end{equation}
Primary metric: AUROC; secondary: F1-macro and balanced accuracy.

\subsection{Phase~2: Random survival forest}\label{sec:rsf}

A random survival forest (RSF;~\cite{[CITE:Ishwaran2008]};
$n_\text{estimators}=100$, \texttt{random\_state}=42) was trained on
the same temporal splits as the classification experiments; a single
survival decision tree (survival-tree analogue of the classification DT)
was trained under the same protocol as its tree-based counterpart.
EFS was defined as time from diagnosis to first of: relapse, induction
failure, or death; OS was time to death.
Censoring was at last follow-up.
Evaluation used Harrell's C-index (Eq.~\ref{eq:cindex}), IPCW
time-dependent AUC (Eq.~\ref{eq:itauc}) at $t^*\in\{12,24,36\}$
months, and the log-rank test on Kaplan-Meier curves dichotomised at
the median risk score.

\subsection{External OOD evaluation}\label{sec:ood}

For external validation, all $n=301$ internal patients served as the
training set and external patients served as the test set; no
site-specific fine-tuning was performed.
DT and RF were trained on the full internal set using the same
preprocessing pipeline (Eq.~\ref{eq:zscore}).
RSF was trained identically.
For Phase~1 LLM evaluation, no training was involved; the same
zero-shot prompts were applied to external patient feature vectors.

\subsection{RAG belief revision and LLM narration}\label{sec:rag_methods}

For test patient $i$, the $K=5$ most similar training cases were
retrieved by cosine similarity in the standardised feature space:
\begin{equation}
\mathcal{N}(i) = \underset{j \in \mathcal{D}^{\text{train}}}
  {\arg\text{top-}K}\
  \frac{\mathbf{z}_i \cdot \mathbf{z}_j}
       {\|\mathbf{z}_i\|_2 \|\mathbf{z}_j\|_2}, \quad K = 5.
\label{eq:rag}
\end{equation}
The RAG evidence probability is
\begin{equation}
p_{\text{rag}} =
  \frac{\sum_{j \in \mathcal{N}(i)} s_j\, y_j}
       {\sum_{j \in \mathcal{N}(i)} s_j},
\label{eq:prag}
\end{equation}
where $s_j$ is the cosine similarity to neighbour $j$ and
$y_j \in \{0,1\}$ is the binary task outcome (or event indicator
for survival).
The revised probability is
\begin{equation}
p_{\text{rev}} = \alpha\, p_0 + (1-\alpha)\, p_{\text{rag}},
\qquad
\alpha = 1 - (1-\alpha_{\text{base}})\sqrt{\bar{s}},
\label{eq:prev}
\end{equation}
where $\alpha_\text{base}=0.6$ and
$\bar{s} = \frac{1}{K}\sum_{j \in \mathcal{N}(i)} s_j$.
For survival, $p_0$ is the min-max normalised RSF risk score.

After belief revision, an LLM prompt containing the patient's features,
DT decision path, top-6 salient features (ranked by
$\text{importance}_j \times |z_{ij}|$), and top-5 retrieved case
profiles was submitted to Gemini-3.1-Flash-Lite to generate a 3--5
sentence clinical narrative.

\subsection{Survival metrics}\label{sec:surv_methods}

Harrell's concordance index:
\begin{equation}
\hat{C} =
\frac{
  \sum_{i \neq j}
  \mathbf{1}[T_i < T_j]\,
  \mathbf{1}[\delta_i = 1]\,
  \mathbf{1}[\hat{r}_i > \hat{r}_j]
}{
  \sum_{i \neq j}
  \mathbf{1}[T_i < T_j]\,
  \mathbf{1}[\delta_i = 1]
}.
\label{eq:cindex}
\end{equation}
IPCW time-dependent AUC:
\begin{equation}
\widehat{\text{AUC}}(t^*) =
\frac{
  \sum_{i:\,T_i \leq t^*,\,\delta_i=1}
  \sum_{j:\,T_j > t^*}
  \mathbf{1}[\hat{r}_i > \hat{r}_j]\,\hat{w}_i
}{
  \sum_{i:\,T_i \leq t^*,\,\delta_i=1}
  \sum_{j:\,T_j > t^*}
  \hat{w}_i
},
\label{eq:itauc}
\end{equation}
where $\hat{w}_i = 1/\hat{G}(T_i)^2$ and $\hat{G}(\cdot)$ is the
Kaplan-Meier censoring estimator~\cite{[CITE:Uno2011]}.

\subsection{Statistical analysis}

AUROC was computed with \texttt{scikit-learn}~\cite{[CITE:sklearn]};
C-index and IPCW AUC with
\texttt{scikit-survival}~\cite{[CITE:sksurv]}.
Bootstrap 95\% confidence intervals (1000 resamples) are reported for
Split~3 ($n_\text{test}=17$).
No multiple-testing correction was applied given the exploratory nature
of this study; all reported $p$-values are two-sided.

%%==============================================================%%
\section*{Data availability}
%%==============================================================%%

The clinical dataset contains patient health information and cannot be
made publicly available due to privacy regulations.
De-identified data may be available through a data access agreement
with [HOSPITAL NAME]; requests to the corresponding author.

%%==============================================================%%
\section*{Code availability}
%%==============================================================%%

Analysis code, LLM prompts, and the bone marrow morphology pipeline
will be made available at
\url{https://github.com/YanbeiJiang/ALL_prediction_NCS}
upon acceptance.

%%==============================================================%%
\section*{Acknowledgements}
%%==============================================================%%

[To be completed — funding sources, clinical collaborators]

%%==============================================================%%
\section*{Author contributions}
%%==============================================================%%

[To be completed]

%%==============================================================%%
\section*{Competing interests}
%%==============================================================%%

The authors declare no competing interests.

%%==============================================================%%
%% Main tables
%%==============================================================%%

\begin{table}[ht]
\caption{Feature set design aligned with the clinical data-acquisition
timeline. All six prediction targets are evaluated at every feature
level; $\dagger$~at $\mathcal{F}_7$, D19MRD and
D46MRD are used as \emph{input features}, so only Final Risk and
Relapse are predicted. EFS/OS are evaluated at all levels using
the same feature vector as classification tasks.}
\label{tab:expdesign}
\centering
\small
\begin{tabular}{lllc}
\toprule
\textbf{Level} & \textbf{Timepoint} & \textbf{Cumulative features added} & \textbf{Class. targets} \\
\midrule
$\mathcal{F}_1$     & Day 0     & WBC count, Age & 4 \\
$\mathcal{F}_2$     & Day 1     & + Blast\% & 4 \\
$\mathcal{F}_3$     & Day 1     & + Immunophenotype + 51 flow cytometry markers & 4 \\
$\mathcal{F}_4$     & Day 1     & $\mathcal{F}_3$ + non-risk clinical factors & 4 \\
$\mathcal{F}_5$     & Day 5     & $\mathcal{F}_3$ + CNS status + Day~5 blast\% + Initial Risk & 4 \\
$\mathcal{F}_6$     & Day 7--10 & $\mathcal{F}_5$ + Fusion gene + Karyotype + FISH & 4 \\
$\mathcal{F}_7$     & Day 46    & $\mathcal{F}_6$ + D19MRD + D46MRD (as inputs) & 2$^\dagger$ \\
\bottomrule
\multicolumn{4}{l}{Survival targets (EFS, OS) evaluated at all levels using the same feature vector.}
\end{tabular}
\end{table}

\begin{table}[ht]
\caption{\textbf{Complete variable specification for each feature level}
(Supplementary Table). Each level adds the listed variables to its base
set; $\mathcal{F}_4$ branches from $\mathcal{F}_3$,
while the main clinical chain is
$\mathcal{F}_1\!\subseteq\!\mathcal{F}_2\!\subseteq\!\mathcal{F}_3\!
\subseteq\!\mathcal{F}_5\!\subseteq\!\mathcal{F}_6\!\subseteq\!\mathcal{F}_7$.}
\label{tab:featurelist}
\centering
\footnotesize
\setlength{\tabcolsep}{4pt}
\begin{tabular}{@{}l p{0.80\linewidth}@{}}
\toprule
\textbf{Level} & \textbf{Variables added (base set)} \\
\midrule
$\mathcal{F}_1$ & WBC count; Age \\
$\mathcal{F}_2$ & $+$ Blast\% \\
$\mathcal{F}_3$ & $+$ Immunophenotype (B/T lineage); 42 CD markers (CD1a, CD2, CD3, cCD3, CD4, CD5, CD7, CD8, CD9, CD10, CD11b, CD11c, CD13, CD14, CD15, CD19, CD20, CD22, cCD22, CD24, CD33, CD34, CD38, CD41, CD44, CD45, CD58, CD61, CD64, CD65, CD66c, CD71, CD73, cCD79a, CD86, CD90, CD97, CD99, CD117, CD123, CD133, CD200); 9 additional markers (GPA, cMPO, NG2, TdT, c$\mu$, sIgM, TCR$\alpha\beta$, TCR$\gamma\delta$, HLA-DR) \\
$\mathcal{F}_4$ & (base $\mathcal{F}_3$) $+$ non-risk clinical factors: chief-complaint category, pregnancy number, birth order, term/preterm status, delivery mode, birth weight, feeding type, admission weight, liver span, spleen span, haemoglobin, platelet count \\
$\mathcal{F}_5$ & (base $\mathcal{F}_3$) $+$ CNS status (CNSL); Day-5 peripheral blast\%; Initial Risk \\
$\mathcal{F}_6$ & (base $\mathcal{F}_5$) $+$ Fusion gene; Karyotype; FISH \\
$\mathcal{F}_7$ & (base $\mathcal{F}_6$) $+$ D19MRD; D46MRD (as input features) \\
$\mathcal{F}_k^{+\text{img/morph/both}}$ & (base $\mathcal{F}_k$, $k\in\{3,\dots,7\}$) $+$ raw bone marrow smear images ($+$img), the structured morphology profile from the SAM/BiomedCLIP pipeline: per-patient blast-cell descriptors (area, perimeter, circularity, eccentricity, aspect ratio, nucleus-to-cytoplasm ratio, GLCM texture, chromatin fineness, nucleoli count/size, vacuole \%, hand-mirror \%) and cell-type proportions ($+$morph), or both ($+$both) \\
\bottomrule
\end{tabular}
\end{table}

%%==============================================================%%
%% Figure captions
%%==============================================================%%

\begin{figure}[ht]
\caption{\textbf{Overview of the temporally-structured multimodal
prediction framework.}
\textbf{a,}~The three parallel timelines: feature acquisition
(Day~0--46, $\mathcal{F}_1$–$\mathcal{F}_{7}$), prediction targets
(six tasks: Final Risk, D19MRD, D46MRD, Relapse, EFS, OS), and the
dynamic case database (patient enrolment year 2015--2019, temporal
splits).
\textbf{b,}~Phase~1 architecture: four LLMs receive serialised patient
features and output class labels or risk scores without task-specific
training.
\textbf{c,}~Phase~2 architecture: a decision tree trained on historical
cases generates predictions and decision paths; RAG retrieves the
$K=5$ most similar cases; LLM synthesises a clinical narrative.
\textbf{d,}~Evaluation scope: internal temporal splits (three splits) and
external OOD validation (merged external cohort).}
\label{fig:overview}
\end{figure}

\begin{figure}[ht]
\caption{\textbf{Phase~1 feature-level ablation across six prediction
targets.}
\textbf{a,}~Final Risk AUROC vs feature level for four LLMs.
\textbf{b,}~Relapse AUROC. \textbf{c,}~D19MRD AUROC.
\textbf{d,}~D46MRD AUROC. \textbf{e,}~EFS C-index.
\textbf{f,}~OS C-index.
Vertical dashed lines indicate clinical timepoints: Day~1 (flow
cytometry), Day~5 (early response), Day~7--10
(cytogenetics), Day~46 (MRD, $\mathcal{F}_7$).
Open symbols: internal cohort. Filled symbols: External cohort.}
\label{fig:ablation}
\end{figure}

\begin{figure}[ht]
\caption{\textbf{Bone marrow morphology pipeline and its predictive
contribution.}
\textbf{a,}~Pipeline schematic: SAM cell segmentation
$\rightarrow$ morphological feature extraction $\rightarrow$
BiomedCLIP cell-type clustering $\rightarrow$ per-patient morphology
profile.
\textbf{b,}~Representative segmented cells with BiomedCLIP cell-type
assignments.
\textbf{c,}~Predictive contribution of the morphology features in
Phase~2: base versus $+$morph RF AUROC for D19MRD at $\mathcal{F}_3$;
the systematic base-vs-$+$morph comparison across all feature levels
($\mathcal{F}_3$--$\mathcal{F}_7$, RF/RSF, internal and external) is
in Table~\ref{tab:phase2_morph}.}
\label{fig:morph}
\end{figure}

\begin{figure}[ht]
\caption{\textbf{Phase~2 temporal learning curves and classification
results.}
\textbf{a,}~Final Risk DT AUROC across feature levels for three
temporal splits; shaded band: 95\% bootstrap CI for Split~3
($n_\text{test}=17$).
\textbf{b,}~Training-set size vs AUROC for Final Risk and Relapse
(DT and RF) at $\mathcal{F}_7$.
\textbf{c,}~External cohort RF AUROC comparison with internal
Split~2 RF AUROC across all four classification tasks.}
\label{fig:p2ablation}
\end{figure}

\begin{figure}[ht]
\caption{\textbf{Survival prediction: Kaplan-Meier curves and
C-index summary.}
\textbf{a,}~EFS Kaplan-Meier curves (high vs low risk, median split)
for the best Phase~1 model (Gemma, $\mathcal{F}_7$, internal cohort).
\textbf{b,}~EFS KM curves for Phase~2 RSF (Split~2, $\mathcal{F}_7$,
log-rank $p=0.010$).
\textbf{c,}~OS KM curves for Phase~2 RSF on the external cohort
(log-rank $p=0.031$).
\textbf{d,}~C-index heatmap: all LLMs $\times$ all feature levels
for EFS (internal, left; external cohort, right).
Numbers at risk and 95\% CI shown.}
\label{fig:km}
\end{figure}

\begin{figure}[ht]
\caption{\textbf{RAG belief revision and representative case narratives.}
\textbf{a,}~AUROC before vs after belief revision across four tasks
and three splits; point size proportional to training-set size.
Filled points: internal; open points: External cohort.
\textbf{b--c,}~Two representative cases from $\mathcal{F}_7$/Split~3.
For each case: patient feature summary, DT decision path, top-3
retrieved similar historical cases with outcomes, and LLM-generated
clinical narrative.
Case~b: correctly classified IR patient with persistent MRD and
concordant historical neighbours.
Case~c: borderline case with conflicting MRD clearance vs cytogenetic
risk signals.}
\label{fig:rag}
\end{figure}

%%==============================================================%%
%% References
%%==============================================================%%
\begin{thebibliography}{99}

\bibitem{[CITE:Hunger2015]}
Hunger, S.P. \& Mullighan, C.G.
Acute lymphoblastic leukemia in children.
\textit{N. Engl. J. Med.} \textbf{373}, 1541--1552 (2015).

\bibitem{[CITE:Bhojwani2013]}
Bhojwani, D. \& Pui, C.-H.
Relapsed childhood acute lymphoblastic leukaemia.
\textit{Lancet Oncol.} \textbf{14}, e205--e217 (2013).

\bibitem{[CITE:Pui2008]}
Pui, C.-H., Robison, L.L. \& Look, A.T.
Acute lymphoblastic leukaemia.
\textit{Lancet} \textbf{371}, 1030--1043 (2008).

\bibitem{[CITE:Good2018]}
Good, Z. et al.
Single-cell developmental classification of B cell precursor acute
lymphoblastic leukemia at diagnosis reveals predictors of relapse.
\textit{Nat. Med.} \textbf{24}, 474--483 (2018).

\bibitem{[CITE:Fitter2021]}
Fitter, S. et al.
CKLF and IL1B transcript levels at diagnosis are predictive of
relapse in children with pre-B-cell acute lymphoblastic leukaemia.
\textit{Br. J. Haematol.} \textbf{193}, 107--118 (2021).

\bibitem{[CITE:Mosquera2024]}
Mosquera Orgueira, A. et al.
Refining risk prediction in pediatric acute lymphoblastic leukemia
through DNA methylation profiling.
\textit{Clin. Epigenetics} \textbf{16}, 53 (2024).

\bibitem{[CITE:Pan2017]}
Pan, L. et al.
Machine learning applications for prediction of relapse in childhood
acute lymphoblastic leukemia.
\textit{Sci. Rep.} \textbf{7}, 7402 (2017).

\bibitem{[CITE:Martinez2025]}
Mart\'inez-Rubio, \'{A}. et al.
Computational flow cytometry immunophenotyping at diagnosis is unable
to predict relapse in childhood B-cell acute lymphoblastic leukemia.
\textit{Comput. Biol. Med.} \textbf{185}, 109831 (2025).

\bibitem{[CITE:Steyerberg2019]}
Steyerberg, E.W.
\textit{Clinical Prediction Models} 2nd edn (Springer, 2019).

\bibitem{[CITE:Phaterpekar2026]}
Phaterpekar, T. et al.
Investigating fine-tuning versus zero-shot learning for general large
language models when predicting cancer survival from initial oncology
consultation documents.
\textit{ESMO Real World Data Digit. Oncol.} (2026).
doi:10.1016/j.esmorw.2026.100703.

\bibitem{[CITE:Conter2010]}
Conter, V. et al.
Molecular response to treatment redefines all prognostic factors in
children and adolescents with B-cell precursor acute lymphoblastic
leukemia: results in 3184 patients of the AIEOP-BFM ALL 2000 study.
\textit{Blood} \textbf{115}, 3206--3214 (2010).

\bibitem{[CITE:Rudin2019]}
Rudin, C.
Stop explaining black box machine learning models for high stakes
decisions and use interpretable models instead.
\textit{Nat. Mach. Intell.} \textbf{1}, 206--215 (2019).

\bibitem{[CITE:SAM]}
Ravi, N. et al.
SAM~2: Segment anything in images and videos.
\textit{arXiv} 2408.00714 (2024).

\bibitem{[CITE:BiomedCLIP]}
Zhang, S. et al.
BiomedCLIP: a multimodal biomedical foundation model pretrained from
fifteen million scientific image-text pairs.
\textit{arXiv} 2303.00915 (2023).

\bibitem{[CITE:vLLM]}
Kwon, W. et al.
Efficient memory management for large language model serving with
PagedAttention.
\textit{Proc. ACM SOSP} (2023).

\bibitem{[CITE:Uno2011]}
Uno, H. et al.
On the C-statistics for evaluating overall adequacy of risk prediction
procedures with censored survival data.
\textit{Stat. Med.} \textbf{30}, 1105--1117 (2011).

\bibitem{[CITE:Ishwaran2008]}
Ishwaran, H. et al.
Random survival forests.
\textit{Ann. Appl. Stat.} \textbf{2}, 841--860 (2008).

\bibitem{[CITE:sklearn]}
Pedregosa, F. et al.
Scikit-learn: machine learning in Python.
\textit{J. Mach. Learn. Res.} \textbf{12}, 2825--2830 (2011).

\bibitem{[CITE:sksurv]}
P\"olsterl, S.
scikit-survival: a library for time-to-event analysis built on top
of scikit-learn.
\textit{J. Mach. Learn. Res.} \textbf{21}, 1--6 (2020).

\end{thebibliography}

\end{document}
